\subsection{内乘：余代数情形}

设 \(E = \bigoplus_{p \geq 0} E_p\) 是一个余结合、余单位的分次余代数。那么我们知道（第 2 号，命题 1 和 3），分次对偶 \(E^{*gr}\)（按照本段开头关于分次的约定）具有 A 上的一个 N 型分次代数结构，该代数中两个元素 \(u, v\) 的乘积由 \(uv = m \circ (u \otimes v) \circ c\) 定义，其中 \(c: E \to E \otimes_A E\) 是余积，而 \(m: A \otimes_A A \to A\) 定义乘法。换言之，若对 \(x \in E\) 有 \(c(x) = \sum_i y_i \otimes z_i\) ，则我们可以写成（将 \(A \otimes_A E\) 与 \(E\) 典范地等同）

\[
\begin{array}{r l} \langle x, u v \rangle = (u v) (x) & = \sum_ {i} u (y _ {i}) v (z _ {i}) = v \left(\sum_ {i} u (y _ {i}) z _ {i}\right) \\ & = v (((u \otimes 1 _ {\mathbb {E}}) \circ c) (x)) = \langle ((u \otimes 1 _ {\mathbb {E}}) \circ c) (x), v \rangle . \end{array}
\]

这可以解释为：对 \(E^{*gr}\) 中所有次数为 p 的齐次元素 u，左乘 \(e(u): v \mapsto uv\) 在 \(E^{*gr}\) 中是 E 的次数为 -p 的分次自同态的分次转置

\[
i (u) = (u \otimes 1 _ {\mathrm{E}}) \circ c\tag{46}
\]

因此，在上述记号中，

\[
(i (u)) (x) = \sum_ {i} u (y _ {i}) z _ {i}.
\]

公式 (46) 同样定义了一个元素 \(i(u)\in\mathrm{Endgr}_{\mathbf{A}}(\mathbf{E})\)（对每个元素 \(u\in\mathbf{E}^{*\mathbf{g}\mathbf{r}}\) 而言）；对所有 \(x\in\mathbf{E}\) 及所有 \(u\in\mathbf{E}^{*\mathbf{g}\mathbf{r}}\) ，我们记

\[
x \llcorner u = (i (u)) (x)\tag{47}
\]

使得，对 E*gr 中的 u 和 v，

\[
\langle x, u v \rangle = \langle x \llcorner u, v \rangle .
\]

E 的元素 \(x \llcorner u\) 称为 x 被 u 的右内乘。

这里再次将“作用于”x的元素u放在L中水平线的自由端。

对任意两个元素 \(u, v\)（属于 \(\mathbf{E}^{*\mathrm{gr}}\)），

\[
x \llcorner (u v) = (x \llcorner u) \llcorner v,\tag{48}
\]

换句话说

\[
i (u v) = i (v) \circ i (u).
\]

如上所述，令 \(c(x) = \sum_{i} y_{i} \otimes z_{i}\) ，则有 \(x \llcorner (uv) = \sum_{i} (uv)(y_{i})z_{i}\) 。若

\[
c (y _ {i}) = \sum_ {j} y _ {i j} ^ {\prime} \otimes y _ {i j} ^ {\prime \prime},
\]

那么

\[
x \llcorner (u v) = \sum_ {i, j} u \left(y _ {i j} ^ {\prime}\right) v \left(y _ {i j} ^ {\prime \prime}\right) z _ {i}.\tag{49}
\]

另一方面，如果 \(c(z_{i}) = \sum_{k} z_{ik}' \otimes z_{ik}''\) ，则

\[
(x \llcorner u) \llcorner v = \sum_ {i, k} u (y _ {i}) v (z _ {i k} ^ {\prime}) z _ {i k} ^ {\prime \prime}.\tag{50}
\]

现在，E 的余结合性表明（第 2 号，命题 1）

\[
\sum_ {i, j} y _ {i j} ^ {\prime} \otimes y _ {i j} ^ {\prime \prime} \otimes z _ {i} = \sum_ {i, k} y _ {i} \otimes z _ {i k} ^ {\prime} \otimes z _ {i k} ^ {\prime \prime}\tag{51}
\]

而表达式 (49) 与 (50) 的相等性来自如下事实：它们分别是 (51) 式左、右两端在从 \(E \otimes E \otimes E\) 到 E 的线性映射 f 下的像，其中 \(f(x \otimes y \otimes z) = u(x)v(y)z\) 。

另一方面，我们回顾（第 2 号，命题 3），代数 \(\mathbf{E}^{*\mathrm{gr}}\) 的单位元是线性形式 \(e: x \mapsto \gamma(x).1\) ；因此

\[
x \llcorner e = \sum_ {i} \gamma (y _ {i}) z _ {i} = x
\]

这是余单位定义的直接推论。由于映射 \(u \mapsto i(u)\) 是线性的，可以看出外合成律 \((u, x) \mapsto x \llcorner u\) 在 E 上定义了一个右 E*gr-模结构。

类似地，我们对所有 \(u \in \mathrm{E}^{*\mathbf{g}\mathbf{r}}\) 定义 \(\mathbf{E}\) 的自同态

\[
i ^ {\prime} (u) = \left(1 _ {E} \otimes u\right) \circ c\tag{52}
\]

并且，对所有 \(x \in \mathbf{E}\) ，我们记

\[
\left(i ^ {\prime} (u)\right) (x) = u \lrcorner x\tag{53}
\]

并且 E 的这个元素称为 x 被 u 的左内乘。如前所述可以看出，外合成律 \((u, v) \mapsto u \lrcorner x\) 在 E 上定义了一个左 \(\mathrm{E}^{*gr}\)-模结构。此外，这两个结构是相容的，换言之，

\[
(u \lrcorner x) \llcorner v = u \lrcorner (x \llcorner v)\tag{54}
\]

这对 E*gr 中的 \(u, v\)（II，§ 1，第 14 号）成立。采用与上述相同的记号，(54) 的左边是 \(\sum_{i,j} u(z_i) v(y_{ij}') y_{ij}''\) ，右边是 \(\sum_{i,k} v(y_i) u(z_{ik}') z_{ik}''\) ；它们的相等性来自如下事实：它们分别是 (51) 式左、右两端在从 E ⊗ E ⊗ E 到 E 的线性映射 \(g\) 下的像，其中 \(g(x \otimes y \otimes z) = v(x) u(z)y\) 。

因此可以看出，E 上的两个外合成律在此集合上定义了一个 \((\mathrm{E}^{*}\mathrm{gr}, \mathrm{E}^{*}\mathrm{gr})\)-双模结构。

当余代数 E 余交换时，则 \(u \lrcorner x = x \llcorner u\) ，其中 \(x \in \mathbf{E}\) 和 \(u \in \mathrm{E}^{*\mathrm{gr}}\) 任取；当它余反交换（§ 4，第 9 号）且 \(u \in \mathrm{E}_p^*\) 、 \(x \in \mathrm{E}_q\) 时，我们可以写成 \(c(x) = \sum_{0 \leqslant j \leqslant q} \left( \sum_i y_{ij} \otimes z_{i,q-j} \right)\) ，其中 \(y_{ij}\) 和 \(z_{ij}\) 都属于 \(\mathrm{E}_j\) （对每个 \(j\) ），于是由假设

\[
\sum_ {i} z _ {i j} \otimes y _ {i, q - j} = (- 1) ^ {j (q - j)} \sum_ {i} y _ {i j} \otimes z _ {i, q - j}.
\]

由定义， \(x \llcorner u = \sum_{0 \leqslant j \leqslant q} \left( \sum_{i} u(y_{ij}) z_{i,q-j} \right)\) 和

\[
u \lrcorner x = \sum_ {0 \leqslant j \leqslant q} \left(\sum_ {i} u (z _ {i, q - j}) y _ {i j}\right).
\]

由于 \(u(y_{ij}) = 0\) （相应地 \(u(z_{i,q-j}) = 0\) ），除非 \(j = p\) （相应地 \(q - j = p\) ），由上述可知 \(u \lrcorner x = (-1)^{p(q-p)}x \llcorner u\) 。

最后，设 \(\phi: \mathrm{E} \to \mathrm{F}\) 是一个分次余代数态射；那么已经看到（第 2 号），分次转置 \(^t\phi: \mathrm{F}^{*\mathrm{gr}} \to \mathrm{E}^{*\mathrm{gr}}\) 是一个分次代数同态；因此，对 \(x \in \mathrm{E}, u, v\) （属于 \(\mathrm{F}^{*\mathrm{gr}}\) ），

\[
\begin{array}{r l} \langle \phi (x \llcorner {} ^ {t} \phi (u)), v \rangle & = \langle x \llcorner {} ^ {t} \phi (u), ^ {t} \phi (v) \rangle = \langle x, ^ {t} \phi (u) ^ {t} \phi (v) \rangle = \langle x, ^ {t} \phi (u v) \rangle \\ & = \langle \phi (x), u v \rangle = \langle \phi (x) \llcorner u, v \rangle \end{array}
\]

从而

\[
\phi (x) \llcorner u = \phi (x \llcorner {} ^ {t} \phi (u));\tag{55}
\]

并且类似地

\[
u \lrcorner \phi (x) = \phi (^ {t} \phi (u) \lrcorner x).\tag{56}
\]

换言之， \(\phi\) 是一个 \(\mathbf{F}^{*\mathrm{gr}}\) -同态，从 \((\mathrm{E}^{*\mathrm{gr}}, \mathrm{E}^{*\mathrm{gr}})\) -双模 E 到 \((\mathrm{F}^{*\mathrm{gr}}, \mathrm{F}^{*\mathrm{gr}})\) -双模 F，相对于环同态

\[
{ } ^ { t } \phi : \mathrm{F} ^ { * \mathrm{gr} } \rightarrow \mathrm{E} ^ { * \mathrm{gr} } .
\]

例子。特别地，当 E 是分次余代数 \(\mathsf{T}(\mathbf{M})\) , \(\mathsf{S}(\mathbf{M})\) 或 \(\bigwedge(\mathbf{M})\) 之一（M 为一个 A-模）（第 1 号，例 5、6 和 7）时，可以应用上述内容。为了显式求出如此得到的双模结构，我们再次将 \(\mathsf{T}(\mathbf{M})^{*\mathbf{g}\mathbf{r}}\) （相应地 \(\mathsf{S}(\mathbf{M})^{*\mathbf{g}\mathbf{r}}\) ，相应地 \(\bigwedge(\mathbf{M})^{*\mathbf{g}\mathbf{r}}\) ）中次数为 n 的齐次元素 f 等同于 \(M^{n}\) 上的一个 n-线性形式（相应地，对称 n-线性形式；相应地，交错 n-线性形式，也称为 n-形式）。只需就 M 的元素表出这些乘积

\[
\left(x _ {1} \otimes x _ {2} \otimes \dots \otimes x _ {p}\right) \llcorner f (\text { resp. } (x _ {1} x _ {2} \dots x _ {p}) \llcorner f,
\]

（相应地 \((x_{1} \wedge x_{2} \wedge \cdots \wedge x_{p}) \llcorner f)\) ），对 M 中元素的每一有限序列 \((x_{i})_{1 \leqslant i \leqslant p}\) ，以及左内乘的类似表达式。现在，定义 (46) 和 (52) 以及第 1 号的公式 (3)、(6) 和 (9) 分别给出：

\[
\left\{ \begin{array}{l} (x _ {1} \otimes x _ {2} \otimes \dots \otimes x _ {p}) \llcorner f = f \lrcorner (x _ {1} \otimes x _ {2} \otimes \dots \otimes x _ {p}) = 0 \\ (x _ {1} x _ {2} \ldots x _ {p}) \llcorner f = f \lrcorner (x _ {1} x _ {2} \ldots x _ {p}) = 0 \quad \text {for} \quad p <   n. \\ (x _ {1} \wedge x _ {2} \wedge \dots \wedge x _ {p}) \llcorner f = f \lrcorner (x _ {1} \wedge x _ {2} \wedge \dots \wedge x _ {p}) = 0 \end{array} \right.\tag{57}
\]

对 \(p \geqslant n\) ，我们分别有

\begin{enumerate}
\def\labelenumi{(\arabic{enumi})}
\setcounter{enumi}{57}
\tightlist
\item
\end{enumerate}

\[
\left(x _ {1} \otimes x _ {2} \otimes \dots \otimes x _ {p}\right) \llcorner f = f \left(x _ {1}, \dots , x _ {n}\right) x _ {n + 1} \otimes \dots \otimes x _ {p}\tag{59}
\]

\[
\left(x _ {1} x _ {2} \dots x _ {p}\right) \llcorner f = \sum_ {\sigma} f \left(x _ {\sigma (1)}, \dots , x _ {\sigma (n)}\right) x _ {\sigma (n + 1)} \dots x _ {\sigma (p)}\tag{60}
\]

\[
\left(x _ {1} \wedge x _ {2} \wedge \dots \wedge x _ {p}\right) \llcorner f = \sum_ {\sigma} \varepsilon_ {\sigma} f \left(x _ {\sigma (1)}, \dots , x _ {\sigma (n)}\right) x _ {\sigma (n + 1)} \wedge \dots \wedge x _ {\sigma (p)}
\]

（其中，在 (59) 和 (60) 中，求和是对排列 \(\sigma \in \mathfrak{S}_p\) 中在每个区间 \([1, n]\) 和 \([n + 1, p]\) （均为 \(\mathbf{N}\) 的区间）上递增的那些进行的）；并且类似地

\begin{enumerate}
\def\labelenumi{(\arabic{enumi})}
\setcounter{enumi}{60}
\tightlist
\item
\end{enumerate}

\[
f \lrcorner (x _ {1} \otimes x _ {2} \otimes \dots \otimes x _ {p}) = f (x _ {p - n + 1}, \dots , x _ {p}) x _ {1} \otimes x _ {2} \otimes \dots \otimes x _ {p - n}\tag{62}
\]

\[
f \lrcorner (x _ {1} x _ {2} \dots x _ {p}) = \sum_ {\sigma} f (x _ {\sigma (p - n + 1)}, \dots , x _ {\sigma (p)}) x _ {\sigma (1)} \dots x _ {\sigma (p - n)}\tag{63}
\]

\[
f \lrcorner (x _ {1} \wedge x _ {2} \wedge \dots \wedge x _ {p}) =
\]

\[
\sum_ {\sigma} \varepsilon_ {\sigma} f (x _ {\sigma (p - n + 1)}, \dots , x _ {\sigma (p)}) x _ {\sigma (1)} \wedge \dots \wedge x _ {\sigma (p - n)}
\]

（其中，在 (62) 和 (63) 中，求和是对排列 \(\sigma \in \mathfrak{S}_p\) 中在每个区间 \([1, p - n]\) 和 \([p - n + 1, p]\) （均为 \(\mathbf{N}\) 的区间）上递增的那些进行的）。

\subsection{内乘：双代数情形}

设 E 是一个分次双代数（相应地，斜分次双代数）（第 4 号，定义 3）；那么第 6 号和第 7 号的结果可以用来定义右（相应地，左）内乘 \(x \llcorner u \in \mathbf{E}\) 和 \(u \llcorner x \in \mathbf{E}^{*\mathbf{gr}}\) （相应地 \(u \lrcorner x \in \mathbf{E}\) 和 \(x \lrcorner u \in \mathbf{E}^{*\mathbf{gr}}\) )

这对所有 \(x \in \mathbf{E}\) 及所有 \(u \in \mathbf{E}^{*\mathbf{gr}}\) 进行。于是就在 E 上得到一个 (E, E)-双模结构和一个 \((\mathbf{E}^{*\mathbf{gr}}, \mathbf{E}^{*\mathbf{gr}})\)-双模结构。进一步：

命题 9. 设 E 是一个分次双代数（相应地，斜分次双代数）。对 E 中次数为 1 的每个元素 x，被 x 所作的左内乘和右内乘都是代数 E*gr 的导子（相应地，反导子）（§ 10，第 2 号）。

在第 6 号的记号中，对分次双代数（相应地，斜分次双代数）E 中每个次数为 1 的齐次元素 \(x\) ，

\[
c (x) = x \otimes 1 + 1 \otimes x,
\]

这由第 3 号命题 5 及 \(c\) 是次数为 0 的同态这一事实得出。首先假设 \(E\) 是分次双代数。对所有 \(y \in E\) ，根据定义

\[
\langle (u v) \llcorner x, y \rangle = \langle u v, x y \rangle = m ((u \otimes v) (c (x y)))
\]

并且由于 \(c\) 是代数同态， \(c(xy) = c(x)c(y)\) 。设 \(c(y) = \sum_{i} s_i \otimes t_i\) ，其中 \(s_i\) 和 \(t_i\) 属于 \(E\) ；因此

\[
c (x y) = \sum_ {i} (x s _ {i}) \otimes t _ {i} + \sum_ {i} s _ {i} \otimes (x t _ {i}).
\]

。因此 \(\langle (uv) \llcorner x, y \rangle = \sum_{i} u(xs_i)v(t_i) + \sum_{i} u(s_i)v(xt_i)\) 。但我们可以写

\[
\sum_ {i} u (x s _ {i}) v (t _ {i}) = m (((u \llcorner x) \otimes v) (c (y))) = \langle (u \llcorner x) v, y \rangle
\]

并且类似地

\[
\sum_ {i} u (s _ {i}) v (x t _ {i}) = m ((u \otimes (v \llcorner x)) (c (y))) = \langle u (v \llcorner x), y \rangle ,
\]

由此，回到用 \(i(x)\) 记内乘的记号，

\[
(i (x)) (u v) = ((i (x)) (u)) v + u ((i (x)) (v))\tag{64}
\]

，这证明了 \(i(x)\) 是 \(\mathbf{E}^{*}\mathbf{gr}\) 中的导子。

现在假设 \(\mathbf{E}\) 是斜分次双代数，并设 \(u\in \mathbf{E}_p^*\) , \(v\in \mathbf{E}_q^*\) 和 \(y\in \mathbf{E}_r\) ；则我们可以写

\[
c (y) = \sum_ {0 \leqslant j \leqslant r} \left(\sum_ {i} s _ {i j} \otimes t _ {i, r - j}\right)
\]

其中 \(s_{ij}\) 和 \(t_{ij}\) 属于 \(\mathbf{E}_j\) ；则由 \(\mathbf{E}^{\mathfrak{g}} \otimes_{\mathbf{A}} \mathbf{E}\) 中乘积的定义，

\[
c (x y) = c (x) c (y) = \sum_ {0 \leqslant j \leqslant r} \left(\sum_ {i} \left(x s _ {i j}\right) \otimes t _ {i, r - j} + (- 1) ^ {j} \sum_ {i} s _ {i j} \otimes \left(x t _ {i, r - j}\right)\right)
\]

从而这一次

\[
\langle (u v) \llcorner x, y \rangle = \sum_ {0 \leqslant j \leqslant r} \left(\sum_ {i} u \left(x s _ {i j}\right) v \left(t _ {i, r - j}\right) + (- 1) ^ {j} \sum_ {i} u \left(s _ {i j}\right) v \left(x t _ {i, r - j}\right)\right).
\]

同样地， \(\sum_{0\leqslant j\leqslant r}\left(\sum_{i}u(xs_{ij})v(t_{i,r - j})\right) = \langle (u\llcorner x)v,y\rangle\) 。另一方面， \(u(s_{ij}) = 0\) ，除非 \(j = -p\) ，因此我们也可以写成

\[
\sum_ {0 \leqslant j \leqslant r} (- 1) ^ {j} \left(\sum_ {i} u (s _ {i j}) v (x t _ {i, r - j})\right) = (- 1) ^ {p} \langle u (v \llcorner x), y \rangle .
\]

因此我们得出结论：

\[
(i (x)) (u v) = ((i (x)) (u)) v + (- 1) ^ {p} u ((i (x)) (v)),\tag{65}
\]

换言之， \(i(x)\) 是 \(\mathrm{E}^{*\mathrm{gr}}\) 中的反导子。关于 E 中次数为 1 的元素 \(x\) 所作的左内乘的断言，可类似地证明。

注记. (1) 设 E 是 A 上的分次双代数，且 N, N', N'' 是三个分次 A-模。设 m 是从 N × N' 到 N'' 的 A-双线性映射；对于 u \ensuremath{\in} Homgr\_A(E, N) 且 v \ensuremath{\in} Homgr\_A(E, N')，设 u.v 表示 E 到 N'' 的分次同态 m ∘ (u ⊗ v) ∘ c。另一方面，设 i(x) 表示由 x \ensuremath{\in} E 在 A-模 Homgr\_A(E, N), Homgr\_A(E, N') 与 Homgr\_A(E, N'\,') 中所作的（右或左）内乘。那么，若 x 的次数为 1，

\[
(i (x)) (u. v) = ((i (x)) (u)). v + u. ((i (x)) (v))
\]

这对所有 \(u \in \operatorname{Homgr}_{\mathbf{A}}(\mathbf{E}, \mathbf{N})\) 和 \(v \in \operatorname{Homgr}_{\mathbf{A}}(\mathbf{E}, \mathbf{N}')\) 成立。

在相同条件下，若 E 是斜分次双代数且 u 是次数为 p 的齐次元素，则

\[
(i (x)) (u. v) = ((i (x)) (u)). v + (- 1) ^ {p} u. ((i (x)) (v)).
\]

其证明与命题 9 中的相同。

\begin{enumerate}
\def\labelenumi{(\arabic{enumi})}
\setcounter{enumi}{1}
\tightlist
\item
  上述证明中的同一论证更一般地证明了：对所有 \(x \in \mathbf{E}\) ，若 \(c(x) = \sum_{j} x_j' \otimes x_{j'}'\) ，则对所有 \(u, v\) （属于 \(\mathbf{E}^{*\mathrm{gr}}\) ），“莱布尼茨公式”
\end{enumerate}

\[
(i (x)) (u v) = \sum_ {j} (i (x _ {j} ^ {\prime})) (u) \cdot (i (x _ {j} ^ {\prime \prime})) (v)
\]

成立。特别地，对分次双代数 E 的每个本原元素，即满足 \(c(x) = x \otimes 1 + 1 \otimes x\) 的元素， \(i(x)\) 是 \(\mathrm{E}^{*\mathrm{gr}}\) 的导子。

命题 10. 设 E 是一个分次双代数（相应地，斜分次双代数）。对 E*gr 中次数为 1 的任意元素 f，左内乘与右内乘都是代数 E 的导子（相应地，反导子）。

设 \(x \in \mathrm{E}_p, y \in \mathrm{E}_p\) ( \(p \geqslant 1, q \geqslant 1\) ）。根据第 3 号命题 5，我们可以写出

\[
c (x) = x \otimes 1 + \sum_ {1 \leqslant j \leqslant p - 1} \left(\sum_ {i} x _ {i j} ^ {\prime} \otimes x _ {i j, p - j} ^ {\prime \prime}\right) + 1 \otimes x
\]

\[
c (y) = y \otimes 1 + \sum_ {1 \leqslant k \leqslant q - 1} \left(\sum_ {i} y _ {i k} ^ {\prime} \otimes y _ {i, q - k} ^ {\prime \prime}\right) + 1 \otimes y
\]

（其中 \(x_{ij}'\) 和 \(x_{ij}''\) 属于 \(\mathrm{E}_j\) ， \(y_{ik}'\) 和 \(y_{ik}''\) 属于 \(\mathrm{E}_k\) 。若 \(\mathrm{E}\) 是分次双代数，则 \(c(xy) = c(x)c(y)\) 的属于 \(\mathrm{E}_1 \otimes \mathrm{E}\) 的分量等于

\[
\sum_ {i} x _ {i, 1} ^ {\prime} \otimes x _ {i, p - 1} ^ {\prime \prime} y + \sum_ {i} y _ {i, 1} ^ {\prime} \otimes x y _ {i, q - 1} ^ {\prime \prime}
\]

因此根据定义

\[
\begin{array}{r l} (x y) \llcorner f & = \sum_ {i} f (x _ {i, 1} ^ {\prime}) x _ {i, p - 1} ^ {\prime \prime} y + \sum_ {i} f (y _ {i, 1} ^ {\prime}) x y _ {i, q - 1} ^ {\prime \prime} \\ & = (x \llcorner f) y + x (y \llcorner f) \end{array}
\]

而由 \(f\) 所作的右内乘是一个导子。另一方面，若 \(E\) 是斜分次双代数，则 \(c(xy)\) 的属于 \(E_1 \otimes E\) 的分量等于

\[
\sum_ {i} x _ {i, 1} ^ {\prime} \otimes x _ {i, p - 1} ^ {\prime \prime} y + (- 1) ^ {p} \sum_ {i} y _ {i, 1} ^ {\prime} \otimes x y _ {i, q - 1} ^ {\prime \prime}
\]

而这一次我们得到

\[
(x y) \llcorner f = (x \llcorner f) y + (- 1) ^ {p} x (y \llcorner f)
\]

这表明 \(i(f)\) 是一个反导子。对由 \(f\) 所作的左内乘，论证类似。

例子。命题 9 和 10 特别适用于分次双代数 \(\mathsf{S}(\mathbf{M})\) 以及斜分次双代数 \(\bigwedge (\mathbf{M})\) 。由 \(\mathsf{S}(\mathbf{M})\) （相应地 \(\mathsf{S}(\mathbf{M})^{*\mathrm{gr}}\) ）中次数为 1 的元素所作的内乘是彼此可交换的导子，因为 \(\mathsf{S}(\mathbf{M})\) （相应地 \(\mathsf{S}(\mathbf{M})^{*\mathrm{gr}}\) ）是交换的。

类似地，由 \(\bigwedge (\mathbf{M})\) （相应地 \(\bigwedge (\mathbf{M})^{*\mathrm{gr}}\) ）中次数为 1 的元素所作的内乘是平方为零的反导子，因为代数 \(\bigwedge (\mathbf{M})\) （相应地 \(\bigwedge (\mathbf{M})^{*\mathrm{gr}}\) ）中次数为 1 的元素的平方为零。

\subsection{\texorpdfstring{\(\mathsf{T}(\mathbf{M})\) 与 \(\mathsf{T}(\mathbf{M}^{*})\)、\(\mathsf{S}(\mathbf{M})\) 与 \(\mathsf{S}(\mathbf{M}^{*})\)、\(\bigwedge(\mathbf{M})\) 与 \(\bigwedge(\mathbf{M}^{*})\) 之间的内乘}{T(M) 与 T(M*)、S(M) 与 S(M*)、wedge(M) 与 wedge(M*) 之间的内乘}}

右内乘在 \(\mathsf{T}(\mathbf{M})\) （相应地 \(\mathsf{S}(\mathbf{M})\) ，相应地 \(\bigwedge (\mathbf{M})\) ）上定义了代数 \(\mathsf{T}(\mathbf{M})^{*\mathfrak{g}\mathfrak{r}}\) （相应地 \(\mathsf{S}(\mathbf{M})^{*\mathfrak{g}\mathfrak{r}}\) ，相应地 \(\bigwedge (\mathbf{M})^{*\mathfrak{g}\mathfrak{r}}\) ）上的一个右模结构（第 7 号，例子）。利用第 5 号的典范同态 \(\theta_{\mathsf{T}}\) （相应地 \(\theta_{\mathsf{S}}\) ，相应地 \(\theta_{\wedge}\) ），我们得到

一个右 \(\mathsf{T}(\mathbf{M}^{*})\) -模结构，定义在 \(\mathsf{T}(\mathbf{M})\) 上

一个右 \(\mathbf{S}(\mathbf{M}^{*})\) -模结构，定义在 \(\mathbf{S}(\mathbf{M})\) 上

一个左 \(\bigwedge (\mathbf{M}^{*})\) -模结构，定义在 \(\bigwedge (\mathbf{M})\) 上。

这些结构中任何一个的外合成律也记为

\[
(z ^ {*}, t) \mapsto i (z ^ {*}). t
\]

（通过滥用语言）；我们也写作 \(t \llcorner z^*\) 以代替 \(i(z^*)\) . \(t\) ，在 \(\mathsf{T}(\mathbf{M})\) 或 \(\mathsf{S}(\mathbf{M})\) 的情形；另一方面，我们写作 \(z^* \lrcorner t\) ，在 \(\bigwedge (\mathbf{M})\) 的情形，并称其为 \(t\) 被 \(z^*\) 所作的左内乘，因为此时我们有一个左 \(\bigwedge (\mathbf{M}^*)\) -模法则。对齐次元素 \(z^*\) （次数为 \(n\) ）和齐次元素 \(t\) （次数为 \(p\) ），有 \(i(z^*)\) . \(t = 0\) （若 \(p < n\) ）；而对 \(x_i \in \mathbf{M} (1 \leqslant i \leqslant p)\) , \(x_j^* \in \mathbf{M}^* (1 \leqslant j \leqslant n)\) 及 \(p \geqslant n\) ，根据第 7 号的公式 (58)、(59) 和 (60)，

\[
\begin{array}{r l} i (x _ {1} ^ {*} \otimes x _ {2} ^ {*} \otimes \dots \otimes x _ {n} ^ {*}) \cdot (x _ {1} \otimes x _ {2} \otimes \dots \otimes x _ {p}) \\ & = \left(\prod_ {j = 1} ^ {n} \langle x _ {j} ^ {*}, x _ {j} \rangle\right) x _ {n + 1} \otimes \dots \otimes x _ {p} \end{array}\tag{66}
\]

\[
i \left(x _ {1} ^ {*} x _ {2} ^ {*} \dots x _ {n} ^ {*}\right). (x _ {1} x _ {2} \dots x _ {p}) = \sum_ {\sigma} \left(\prod_ {j = 1} ^ {n} \langle x _ {j} ^ {*}, x _ {\sigma (j)} \rangle\right) x _ {\sigma (n + 1)} \dots x _ {\sigma (p)}\tag{67}
\]

\[
\begin{array}{l} i \left(x _ {1} ^ {*} \wedge x _ {2} ^ {*} \wedge \dots \wedge x _ {n} ^ {*}\right). \left(x _ {1} \wedge x _ {2} \wedge \dots \wedge x _ {p}\right) \\ = (- 1) ^ {n (n - 1) / 2} \sum_ {\sigma} \varepsilon_ {\sigma} \left(\prod_ {j = 1} ^ {n} \langle x _ {j} ^ {*}, x _ {\sigma (j)} \rangle\right) x _ {\sigma (n + 1)} \wedge \dots \wedge x _ {\sigma (p)} \end{array}\tag{68}
\]

其中，在公式 (67) 和 (68) 中， \(\sigma\) 遍历排列 \(\sigma \in \mathfrak{S}_p\) 中在区间 \([1, n]\) 和 \([n + 1, p]\) 上递增的那些排列。我们也可以在内乘记号下写成

\[
\begin{array}{l l} \langle t \llcorner u ^ {*}, v ^ {*} \rangle = \langle t, \theta_ {\mathsf {T}} (u ^ {*} v ^ {*}) \rangle & \text {for} \quad t \in \mathsf {T} (\mathbf {M}), u ^ {*}, v ^ {*} \quad \text {in} \quad \mathsf {T} (\mathbf {M} ^ {*}) \\ \langle t \llcorner u ^ {*}, v ^ {*} \rangle = \langle t, \theta_ {\mathsf {S}} (u ^ {*} v ^ {*}) \rangle & \text {for} \quad t \in \mathsf {S} (\mathbf {M}), u ^ {*}, v ^ {*} \quad \text {in} \quad \mathsf {S} (\mathbf {M} ^ {*}) \\ \langle v ^ {*}, u ^ {*} \lrcorner t \rangle = \langle \theta_ {\wedge} (u ^ {*} \wedge v ^ {*}), t \rangle & \text {for} \quad t \in \bigwedge (\mathbf {M}), u ^ {*}, v ^ {*} \quad \text {in} \quad \bigwedge (\mathbf {M} ^ {*}). \end{array}
\]

我们把显式求出左内乘的类似公式的任务留给读者，这次使用公式 (61)、(62) 和 (63)。

上述内容可以在把 M 换成其对偶 \(M^{*}\) 的情形下应用；此时必须把 \(M^{*}\) 换成双重对偶 \(M^{**}\) ，例如， \(T(M^{*})\) 于是具有代数 \(T(M^{**})\) 上的一个右模结构。但典范映射 \(c_{M}: M \to M^{**}\) 定义了代数同态 \(T(c_{M}): T(M) \to T(M^{**})\) ，借助它， \(T(M^{*})\) 具有一个右 \(T(M)\) -模结构。类似地， \(S(M^{*})\) （相应地 \(\bigwedge(M^{*})\) ）具有一个右 \(S(M)\) -模（相应地，左 \(\bigwedge(M)\) -模）结构。给出这些模的外合成律的显式公式可通过交换 M 与 \(M^{*}\) 的角色从上文直接得出。注意，对所有 \(x \in M\) , \(i(x)\) 总是分次代数 \(S(M^{*})\) （相应地 \(\bigwedge(M^{*})\) ）的导子（相应地，平方为零的反导子）。

命题 11. 典范同态 \(\theta_{\mathsf{T}}: \mathsf{T}(\mathbf{M}^{*}) \to \mathsf{T}(\mathbf{M})^{*\mathsf{gr}}\) （相应地 \(\theta_{\mathsf{S}}: \mathsf{S}(\mathbf{M}) \to \mathsf{S}(\mathbf{M})^{*\mathsf{gr}}\) ，相应地 \(\theta_{\wedge}: \bigwedge (\mathbf{M}^{*}) \to \bigwedge (\mathbf{M})^{*\mathsf{gr}}\) ）是一个右 \(\mathsf{T}(\mathbf{M})\) -模（相应地，右 \(\mathsf{S}(\mathbf{M})\) -模，相应地，左 \(\bigwedge (\mathbf{M})\) -模）同态。

我们首先证明，对 \(z^{*} \in \mathsf{T}(\mathbf{M}^{*})\) 和 \(t \in \mathsf{T}(\mathbf{M})\) ，

\[
\theta_ {T} (z ^ {*} \llcorner t) = \theta_ {T} (z ^ {*}) \llcorner t.\tag{69}
\]

由于 M 是代数 \(\mathsf{T}(\mathbf{M})\) 的生成系，我们只需在 \(t = x \in M\) 时证明 (69)；此外，我们可以把注意力限制在 \(z^{*} = x_{1}^{*} \otimes x_{2}^{*} \otimes \cdots \otimes x_{p}^{*}\) （其中 \(x_{j}^{*} \in M^{*}\) ）的情形；于是，由 (66)（交换 M 与 \(M^{*}\) 的角色）， \(z^{*} \llcorner x = \langle x, x_{1}^{*} \rangle x_{2}^{*} \otimes \cdots \otimes x_{p}^{*}\) 。因此，对 M 中所有 \(y_{2}, \ldots, y_{p}\) ，

\[
\begin{array}{r l} \langle \theta_ {T} (z ^ {*} \llcorner x), y _ {2} \otimes y _ {3} \otimes \dots \otimes y _ {p} \rangle & = \langle x, x _ {1} ^ {*} \rangle \prod_ {j = 2} ^ {p} \langle y _ {j}, x _ {j} ^ {*} \rangle \\ & = \langle \theta_ {T} (z ^ {*}), x \otimes y _ {2} \otimes \dots \otimes y _ {p} \rangle \\ & = \langle \theta_ {T} (z ^ {*}) \llcorner x, y _ {2} \otimes \dots \otimes y _ {p} \rangle \end{array}
\]

由此即得 (69)。

其次我们证明对 \(z^{*} \in \mathbf{S}(\mathbf{M}^{*})\) 和 \(t \in \mathbf{S}(\mathbf{M})\) ,

\[
\theta_ {S} (z ^ {*} \llcorner t) = \theta_ {S} (z ^ {*}) \llcorner t.\tag{70}
\]

如上所述，我们可以限于 \(t = x \in \mathbf{M}\) 的情形。但进一步地，这里 \(i(x)\) 是 \(\mathbf{S}(\mathbf{M}^*)\) 的导子，也是 \(\mathbf{S}(\mathbf{M})^{*\text{gr}}\) 的导子。因此（§ 10，第 7 号，命题 9 的推论），只需对 \(z^* = x^* \in \mathbf{M}^*\) 验证 (70)，因为 \(\mathbf{M}^*\) 是 \(\mathbf{S}(\mathbf{M}^*)\) 的生成系；但这是平凡的，此时两边都等于 \(\langle x^*, x \rangle\) 。类似的论证可证明关系式

\[
\theta_ {\wedge} (t \lrcorner z ^ {*}) = t \lrcorner \theta_ {\wedge} (z ^ {*})\tag{71}
\]

对 \(z^{*} \in \bigwedge(\mathbf{M}^{*})\) 和 \(t \in \bigwedge(\mathbf{M})\) ：此时注意到，对 \(x \in \mathbf{M}\) , \(i(x)\) 在 \(\bigwedge(\mathbf{M}^{*})\) 中以及 \(\bigwedge(\mathbf{M})^{*\text{er}}\) 中都是反导子，并使用 § 10，第 7 号，命题 9 的推论。对左内乘有类似的结果。

\subsection{有限生成自由模情形下内乘的显式形式}

设 \(\mathbf{M}\) 是一个有限生成的自由 A-模， \((e_i)_{1\leq i\leq n}\) 是 \(\mathbf{M}\) 的一组基，而 \((e_i^*)_{1\leq i\leq n}\) 是 \(\mathbf{M}^*\) 的对偶基。对每个有限序列 \(s = (i_1,\ldots ,i_p)\) （由 \([1,n]\) 中元素组成），设 \(e_s = e_{i_1}\otimes e_{i_2}\otimes \dots \otimes e_{i_p}\) （相应地 \(e_s^* = e_{i_1}^*\otimes \dots \otimes e_{i_p}^*\) ）。我们知道（§ 5，第 5 号，定理 1）， \(e_s\) 构成 A-模 \(\mathsf{T}(\mathbf{M})\) 的一组基，而 \(e_s^*\) 构成 A-模 \(\mathsf{T}(\mathbf{M}^*)\) 的一组基。若 \(s,t\) 是两个有限序列（由 \([1,n]\) 中元素组成），设 \(s.t\) 表示如下得到的序列：如果 \(s = (i_1,\ldots ,i_p)\) 和 \(t = (j_1,\ldots ,j_q)\) ，则 \(s.t\) 是序列 \((i_1,\ldots ,i_p,j_1,\ldots ,j_q)\) ，共有 \(p + q\) 项。于是 \(e_{s,t} = e_s\otimes e_t\) ，从而由 (66) 可得

\[
\left\{ \begin{array}{l} e _ {s} \llcorner e _ {t} ^ {*} = 0 \quad \text { if   } s \text {   is   not   of   the   form   } t. u \\ e _ {t. u} \llcorner e _ {t} ^ {*} = e _ {u}. \end{array} \right.\tag{72}
\]

类似地，对称代数 \(\mathbf{S}(\mathbf{M})\) 以单项式 \(e^{\alpha}\) （ \(\alpha \in \mathbf{N}^{n}\) ）的集合为基（§ 6，第 6 号，定理 1），而 \(\mathbf{S}(\mathbf{M}^{*})\) 以单项式 \(e^{*\alpha}\) （ \(\alpha \in \mathbf{N}^{n}\) ）的集合为基；回顾（第 5 号）， \(u_{\alpha}\) （对 \(|\alpha| = k\) ）表示 \((\mathbf{S}^{k}(\mathbf{M}))^{*}\) 的基中与 \((e^{\alpha})_{|\alpha| = k}\) （ \(\mathbf{S}^{k}(\mathbf{M})\) 的基）对偶的元素；因此 \(u_{\alpha}\) （对 \(\alpha \in \mathbf{N}^{n}\) ）构成 \(\mathbf{S}(\mathbf{M})^{*\mathrm{gr}}\) 的一组基。由 \(e^{\beta}\) 所作的右内乘（在 \(\mathbf{S}(\mathbf{M})^{*\mathrm{gr}}\) 中）定义为乘以 \(e^{\beta}\) （在 \(\mathbf{S}(\mathbf{M})\) 中）的映射的转置，由此可知

\[
\left\{ \begin{array}{l l} u _ {\alpha} \llcorner e ^ {\beta} = 0 & \text { if } \quad \alpha \not \geqslant \beta \\ u _ {\alpha} \llcorner e ^ {\beta} = u _ {\alpha - \beta} & \text { if } \quad \alpha \geqslant \beta . \end{array} \right.\tag{73}
\]

类似地，由于 \(\mathbf{S}(\mathbf{M})\) 在此典范地与 \(\mathbf{S}(\mathbf{M})^{*\varepsilon r}\) 的分次对偶等同， \(i(u_{\beta})\) 是乘以 \(u^{\beta}\) （在 \(\mathbf{S}(\mathbf{M})^{*\varepsilon r}\) 中）的分次转置，因此由基 \((u_{\alpha})\) 的乘法表 (33)（第 5 号）推出

\[
\left\{ \begin{array}{l l} e ^ {\alpha} \llcorner u _ {\beta} = 0 & \text { if } \quad \alpha \not \geqslant \beta \\ e ^ {\alpha} \llcorner u _ {\beta} = (\beta , \alpha - \beta) e ^ {\alpha - \beta} & \text { if } \quad \alpha \geqslant \beta . \end{array} \right.\tag{74}
\]

至于 \(\mathsf{S}(\mathbf{M})\) 中元素被 \(\mathsf{S}(\mathbf{M}^*)\) 中元素所作的右内乘，此乘积的定义（第 9 号）以及第 5 号的公式 (34) 使我们能从 (74) 推出以下公式

\[
\left\{ \begin{array}{l l} e ^ {\alpha} \llcorner e ^ {* \beta} = 0 & \text { if } \quad \alpha \not \geqslant \beta \\ e ^ {\alpha} \llcorner e ^ {* \beta} = \frac {\alpha !}{(\alpha - \beta) !} e ^ {\alpha - \beta} & \text { if } \quad \alpha \geqslant \beta . \end{array} \right.\tag{75}
\]

对 \(\mathsf{S}(\mathbf{M}^{*})\) 中元素被 \(\mathsf{S}(\mathbf{M})\) 中元素所作的内乘，存在类似的公式：交换 M 与 \(M^{*}\) 的角色（因为此处 \(M^{**}\) 与 M 等同）。

注记。给定基 \((e_i)_{1\leq i\leq n}\) 使我们能够把代数 \(\mathsf{S}(\mathbf{M})\) 与多项式代数 \(\mathbf{A}[\mathbf{X}_1,\ldots ,\mathbf{X}_n]\) 等同（§ 6，第 6 号）；公式 (75) 表明，由 \(e^{*\alpha}\) 所作的内乘正是微分算子 \(\mathbf{D}^{\alpha} = \mathbf{D}_{1}^{\alpha_{1}}\mathbf{D}_{2}^{\alpha_{2}}\ldots \mathbf{D}_{n}^{\alpha_{n}}\) ，其中 \(\mathbf{D}_i = \partial /\partial \mathbf{X}_i\) （对 \(1\leqslant i\leqslant n\) ）（§ 10，第 11 号，例子）。

最后考虑外代数 \(\bigwedge (\mathbf{M})\) ，它以元素 \(e_{\mathbf{J}}\) 的集合为基，其中 \(\mathbf{J}\) 遍历区间 \([1, n]\) （属于 \(\mathbf{N}\) ）的子集的集合（§ 7，第 8 号，定理 1）；类似地， \(\bigwedge (\mathbf{M}^*)\) 以元素 \(e_{\mathbf{J}}^{*}\) 为基。由第 9 号的公式 (68) 可知

\[
\left\{ \begin{array}{l l} e _ {\mathrm{K}} ^ {*} \lrcorner e _ {\mathrm{J}} = 0 & \text { if } \quad \mathrm{K} \nsubseteq \mathrm{J} \\ e _ {\mathrm{K}} ^ {*} \lrcorner e _ {\mathrm{J}} = (- 1) ^ {p (p - 1) / 2} \rho_ {\mathrm{K}, \mathrm{J} - \mathrm{K}} e _ {\mathrm{J} - \mathrm{K}} & \text { if } \quad \mathrm{K} \subset \mathrm{J} \quad \text { and } \quad p = \operatorname{Card} (\mathrm{K}), \end{array} \right.\tag{76}
\]

（其中 \(\rho_{\mathbf{K},J - \mathbf{K}}\) 是由 § 7 第 8 号公式 (19) 定义的数。把 \(\mathbf{M}\) 和 \(\mathbf{M}^*\) 的角色互换，便有类似的公式。）

\subsection{\texorpdfstring{\(\wedge^{p}(\mathbf{M})\) 与 \(\wedge^{n - p}(\mathbf{M}^{*})\) 之间的同构（\(n\) 维自由模的情形）}{\textbackslash wedge\^{}\{p\}(\textbackslash mathbf\{M\}) 与 \textbackslash wedge\^{}\{n - p\}(\textbackslash mathbf\{M\}\^{}\{*\}) 之间的同构（n 维自由模 M）}}

命题 12. 设 \(\mathbf{M}\) 是一个维数为 \(n\) 的自由 A-模；令 \(e \in \bigwedge^{n}(\mathbf{M})\) 是构成 \(\bigwedge^{n}(\mathbf{M})\) 的一组基的元素，并设 \(e^{*}\) 是 \(\bigwedge^{n}(\mathbf{M}^{*})\) 中使得 \(\{(-1)^{n(n-1)/2} \theta_{\bigwedge}(e^{*})\}\) 是 \(\{e\}\) 在 \((\bigwedge^{n}(\mathbf{M}))^{*}\) 中的对偶基的元素。设 \(\phi: \bigwedge(\mathbf{M}^{*}) \to \bigwedge(\mathbf{M})\) 是映射 \(z \mapsto z \lrcorner e^{*}\) ，而 \(\phi': \bigwedge(\mathbf{M}^{*}) \to \bigwedge(\mathbf{M})\) 是映射 \(z^{*} \mapsto z^{*} \lrcorner e\) 。设 \(\phi_{p}\) （相应地 \(\phi_{p}'\) ）是 \(\phi\) （相应地 \(\phi'\) ）在 \(\bigwedge^{p}(\mathbf{M})\) （相应地 \(\bigwedge^{p}(\mathbf{M}^{*})\) ）上的限制。那么：

\begin{enumerate}
\def\labelenumi{(\roman{enumi})}
\item
  映射 \(\phi\) 是一个左 \(\bigwedge (\mathbf{M})\) -模同构，而映射 \(\phi'\) 是一个左 \(\bigwedge (\mathbf{M}^*)\) -模同构；此外，映射 \(\phi\) 和 \(\phi'\) 互为逆映射。
\item
  映射 \(\phi_p\) 是 A-模 \(\bigwedge^p (\mathbf{M})\) 到 A-模 \(\bigwedge^{n - p}(\mathbf{M}^*)\) 上的同构，而映射 \(\phi_p'\) 是 A-模 \(\bigwedge^p (\mathbf{M}^*)\) 到 A-模 \(\bigwedge^{n - p}(\mathbf{M})\) 上的同构。
\item
  若记 \(\mathbf{B}(u,v^{*}) = \langle u,\theta_{\wedge}(v^{*})\rangle\) （对 \(u\in \bigwedge (\mathbf{M})\) 和 \(v^{*}\in \bigwedge (\mathbf{M}^{*})\) ），则对 \(u^{*}\in \bigwedge^{p}(\mathbf{M}^{*})\) 和 \(v^{*}\in \bigwedge^{n - p}(\mathbf{M}^{*})\)
\end{enumerate}

\[
\mathbf {B} \left(\phi_ {p} ^ {\prime} (u ^ {*}), v ^ {*}\right) = (- 1) ^ {p (n - p)} \mathbf {B} \left(u ^ {*}, \phi_ {n - p} ^ {\prime} (v ^ {*})\right).\tag{77}
\]

\(\phi\) 是 \(\bigwedge (\mathbf{M})\) -线性的以及 \(\phi'\) 是 \(\bigwedge (\mathbf{M}^*)\) -线性的这一事实，由公式 \((u \wedge v) \lrcorner e^* = u \lrcorner (v \lrcorner e^*)\) 和 \((u^* \wedge v^*) \lrcorner e = u^* \lrcorner (v^* \lrcorner e)\) 得出（第 6 号，公式 (37)，利用 \(\theta_{\wedge}\) 是 \(\bigwedge (\mathbf{M}^*)\) 到 \(\bigwedge (\mathbf{M}^*)\) 的反代数上的同构这一事实）。另一方面，存在一个基 \((e_i)_{1 \leq i \leq n}\) （ \(\mathbf{M}\) 的基），使得

\[
e = e _ {1} \wedge e _ {2} \wedge \dots \wedge e _ {n} \quad \text { and } \quad e ^ {*} = (- 1) ^ {n (n - 1) / 2} e _ {1} ^ {*} \wedge e _ {2} ^ {*} \wedge \dots \wedge e _ {n} ^ {*},
\]

（其中 \((e_i^*)\) 是 \((e_i)\) 的对偶基。我们记 \(\mathbf{I} = [1, n]\) ；由 (76) 可知，对每个子集 \(\mathbf{J}\) （属于 \(\mathbf{I}\) ，含 \(p\) 个元素），

\[
\left\{ \begin{array}{l} \phi (e _ {J}) = (- 1) ^ {n (n - 1) / 2 + p (p - 1) / 2} \rho_ {J, I - J} e _ {I - J} ^ {*} \\ \phi^ {\prime} (e _ {J} ^ {*}) = (- 1) ^ {p (p - 1) / 2} \rho_ {J, I - J} e _ {I - J}. \end{array} \right.\tag{78}
\]

这证明了 \(\phi\) 和 \(\phi'\) 是双射；此外， \(\rho_{J,I-J}\rho_{I-J,J}=(-1)^{p(n-p)}\) （ \(\S 7\) ，第 8 号，公式 (21)）；由于数

\[
\frac {n (n - 1)}{2} + \frac {p (p - 1)}{2} + \frac {(n - p) (n - p - 1)}{2} + p (n - p) = n (n - 1)
\]

是偶数，因此得出 \(\phi\) 和 \(\phi'\) 互为逆映射。最后，为证明 (77)，只需取 \(u^{*} = e_{J}^{*}\) 和 \(v^{*} = e_{I - J}^{*}\) ；验证同样由 \(\theta_{\wedge}\) 的定义、公式 (78) 及关系式 \(\rho_{J,I - J}\rho_{I - J,J} = (-1)^{p(n - p)}\) （§ 7，第 8 号，公式 (21)）得出。注意，对 \(u^{*}\in \bigwedge^{p}(\mathbf{M}^{*})\) 和 \(v^{*}\in \bigwedge^{n - p}(\mathbf{M}^{*})\) , \(\mathbf{B}(\phi_{p}^{\prime}(u^{*}), v^{*})\) 是 \(u^{*} \wedge v^{*}\) 关于基 \(\{e^{*}\}\) （ \(\bigwedge^{n}(\mathbf{M}^{*})\) 的基）的系数，至多差一个符号。

命题 13. 在命题 11 的假设和记号下，对 A-模 M 的每个自同态 g，

\[
(\det g) \phi = \bigwedge (^ {t} g) \circ \phi \circ \bigwedge (g).\tag{79}
\]

显然 \(\bigwedge(^{t}g) = \theta_{\wedge}^{-1}\circ (^{t}\bigwedge (g))\circ \theta_{\wedge}\) ；因为 \(\bigwedge (g)\) 是代数 \(\bigwedge (\mathbf{M})\) 的自同态，并且根据定义，对所有

\[
z \in \bigwedge (\mathbf {M}), \quad \theta_ {\wedge} (\bigwedge (g) (z) \lrcorner e ^ {*}) = \theta_ {\wedge} (e ^ {*}) \llcorner \theta_ {\wedge} (\bigwedge (g) (z)),
\]

我们由第 6 号的公式 (42) 推出

\[
\begin{array}{r l} ((\theta_ {\wedge} ^ {- 1} \circ (^ {t} \bigwedge (g)) \circ \theta_ {\wedge}) \circ \phi \circ \bigwedge (g)) (z) & = \theta_ {\wedge} ^ {- 1} (^ {t} \bigwedge (g) (\theta_ {\wedge} (e ^ {*})) \llcorner z) \\ & = z \lrcorner (\bigwedge (t g) (e ^ {*})) = (\det g) (z \lrcorner e ^ {*}) = (\det g) \phi (z) \end{array}
\]

这里用到了 § 8，第 4 号，命题 8。

推论。对 E 的每个自同构 g，

\[
\bigwedge(^ {t} g ^ {- 1}) = (\det g) ^ {- 1} \phi \circ (\bigwedge (g)) \circ \phi^ {- 1}.\tag{80}
\]

\subsection{对与 p-向量相伴的子空间的应用}

设 K 是一个域，E 是 K 上的一个向量空间。回顾一下，对每个 p-向量 \(z \in \bigwedge^{p}(\mathbf{E})\) 相伴着 E 的一个有限维子空间 \(M_{z}\) ，即 E 的使得 \(z \in \bigwedge^{p}(\mathbf{M})\) 的最小向量子空间 M（§ 7，第 2 号，命题 4 的推论）。

命题 14. （i） \(\mathbf{M}_z\) 在 \(\mathbf{E}^*\) 中的正交子空间是 \(x^{*}\in \mathbf{E}^{*}\) （满足 \(x^{*}\lrcorner z = 0\) ）的集合。

\begin{enumerate}
\def\labelenumi{(\roman{enumi})}
\setcounter{enumi}{1}
\tightlist
\item
  子空间 \(\mathbf{M}_z\) （与 \(z\) 相伴）是 \(\bigwedge^{p-1}(\mathbf{E}^*)\) 在映射 \(\lambda_z: u^* \mapsto u^* \lrcorner z\) （从 \(\bigwedge^{p-1}(\mathbf{E}^*)\) 到 \(\mathbf{E}\) ）下的像。
\end{enumerate}

设 \(\mathbf{N}\) 表示 \(\lambda_z\) 的像。对 \(x^* \in \mathbf{E}^*\) 和 \(u^* \in \bigwedge^{p-1}(\mathbf{E}^*)\) ，

\[
\begin{array}{r l} \langle \theta_ {\wedge} (x ^ {*}), u ^ {*} \lrcorner z \rangle = \langle \theta_ {\wedge} (u ^ {*} \wedge x ^ {*}), z \rangle & = (- 1) ^ {p - 1} \langle \theta_ {\wedge} (x ^ {*} \wedge u ^ {*}), z \rangle \\ & = (- 1) ^ {p - 1} \langle \theta_ {\wedge} (u ^ {*}), x ^ {*} \lrcorner z \rangle . \end{array}
\]

因此， \(x^*\) 与 \(\mathbf{N}\) 正交的必要充分条件是 \(x^* \lrcorner z\) 与 \(\theta_\wedge (\bigwedge (\mathrm{E}^*))\) 正交。而后一条件等价于 \(x^* \lrcorner z = 0\) ；因为令 \((e_\lambda)_{\lambda \in L}\) 是 \(\mathbf{E}\) 的一组基；给 \(L\) 一个全序，已知（§ 7，第 8 号，定理 1） \(e_j\) （对 \(J\) 遍历 \(\mathfrak{F}(L)\) （ \(L\) 的有限子集的集合））构成 \(\bigwedge (\mathrm{E})\) 的一组基；于是由第 5 号的公式 (30) 可知，元素 \(\theta_{\wedge}(e_{\mathrm{J}}^{*})\) 是 \(\bigwedge(\mathrm{E})\) 上相对于基 \((e_{\mathrm{J}})\) 的坐标形式（至多差一个符号）；由此得我们的断言。

因此，N 的正交补由 \(x^{*} \in E^{*}\) （满足 \(x^{*} \lrcorner z = 0\) ）组成，从而 (i) 的结论将由 (ii) 得出。

我们首先证明 \(\mathbf{N} \subset \mathbf{M}_z\) 。设 \(\mathbf{M}\) 是 \(\mathbf{E}\) 的一个向量子空间，使得 \(z \in \bigwedge (\mathbf{M})\) ，并设 \(j: \mathbf{M} \to \mathbf{E}\) 是典范单射；令 \(\mu_z\) 表示映射 \(v^* \mapsto v^* \lrcorner z\) （从 \(\bigwedge^{p-1}(\mathbf{M}^*)\) 到 \(\mathbf{M}\) ）；由第 7 号的公式 (60) 可知存在典范分解

\[
\lambda_ {z}: \bigwedge^ {p - 1} (\mathrm{E} ^ {*}) \xrightarrow {\bigwedge^ {p - 1} ({} ^ {t} j)} \bigwedge^ {p - 1} (\mathrm{M} ^ {*}) \xrightarrow {\mu_ {z}} \mathrm{M} \xrightarrow {j} \mathrm{E}
\]

，这证明了 \(\mathbf{N} \subset \mathbf{M}\) ，因此 \(\mathbf{N} \subset \mathbf{M}_z\) （按 \(\mathbf{M}_z\) 的定义）。还需验证 \(\mathbf{N} = \mathbf{M}_z\) 。假设其逆成立：则存在一组基 \((e_i)_{1 \leqslant i \leqslant n}\) （ \(\mathbf{M}_z\) 的基）以及一个元素 \(x^* \in \mathbf{E}^*\) ，使得 \(\langle x^*, e_1 \rangle = 1\) , \(\langle x^*, e_j \rangle = 0\) （对 \(2 \leqslant j \leqslant n\) ），且使得 \(x^*\) 与 \(\mathbf{N}\) 正交，从而 \(x^* \lrcorner z = 0\) 。我们记 \(z = \sum_{\mathbf{H}} a_{\mathbf{H}} e_{\mathbf{H}}\) ，其中求和是对 \([1, n]\) 的含 \(p\) 个元素的子集进行的。由 (68)（第 9 号），

\[
\begin{array}{l l} x ^ {*} \lrcorner e _ {\mathrm{H}} = 0 & \text { if } \quad 1 \notin \mathrm{H} \\ x ^ {*} \lrcorner e _ {\{1 \} \cup \mathrm{H}} = e _ {\mathrm{H}} & \text { if } \quad \mathrm{H} \subset [ 2, n ] \end{array}
\]

这表明关系 \(x^{*} \lrcorner z = 0\) 蕴含 \(a_{\mathbf{H}} = 0\) （对 \(1 \in \mathbf{H}\) ）。但这是不可能的，因为那样 \(z\) 就将属于 \(\bigwedge^{p}(\mathbf{M}')\) ，其中 \(\mathbf{M}'\) 是 \(\mathbf{M}\) 的由 \(e_2, \ldots, e_n\) 生成的子空间。

\subsection{纯 p-向量。格拉斯曼簇}

设 K 是一个域，E 是 K 上的一个向量空间。p-向量 \(z \in \bigwedge^{p}(\mathbf{E})\) 称为纯的（有时也称为可分解的），如果它是非零的并且存在 E 中的向量 \(x_{1}, \ldots, x_{p}\) 使得 \(z = x_{1} \wedge x_{2} \wedge \cdots \wedge x_{p}\) 。为此，必要且充分条件是：与 z 相伴的子空间 \(M_{z}\) （其维数总是 \(\geqslant p\) ，当 \(z \neq 0\) 时）恰好为 p 维（因为此时 \(\bigwedge^{p}(\mathbf{M}_{z})\) 的维数为 1）。特别地，每个非零标量、 \(\mathbf{E} = \bigwedge^{1}(\mathbf{E})\) 中的每个非零元素、以及当 E 的维数为 n 时 \(\bigwedge^{n}(\mathbf{E})\) 中的每个非零元素，都是纯的。

命题 15. 设 E 是一个 n 维向量空间，且设 e 为一个元素 \(\neq0\) （属于 \(\bigwedge^{n}(\mathbf{E})\) ，因而构成此向量空间的一组基）。设 \(\phi:\bigwedge(\mathbf{E})\to\bigwedge(\mathbf{E}^{*})\) 是与 e 相伴的向量空间同构（第 11 号，命题 12）。若 z 是 \(\bigwedge^{p}(\mathbf{E})\) 中的纯元素，则 \(\phi(z)\) 是 \(\bigwedge^{n-p}(\mathbf{E}^{*})\) 中的纯元素，且与 z 和 \(\phi(z)\) 相伴的子空间是正交的。

情形 \(p = 0\) 和 \(p = n\) 是平凡的。因此假设 \(1 \leqslant p \leqslant n - 1\) ，并设 \(z = x_1 \wedge \cdots \wedge x_p \neq 0\) 。于是存在一个基 \((e_i)_{1 \leqslant i \leqslant n}\) （ \(E\) 的），使得 \(e_i = x_i\) （对 \(1 \leqslant i \leqslant p\) ），且 \(e = e_1 \wedge e_2 \wedge \cdots \wedge e_n\) 。于是由第 11 号的公式 (78) 可得 \(\phi(z) = \pm e_{p+1}^* \wedge \cdots \wedge e_n^*\) ，由此得出该命题。

推论. 若 E 的维数为 n，则 E 的每个非零 \((n - 1)\) -向量都是纯的。

命题 16. 元素 \(z \neq 0\) （属于 \(\bigwedge^{p}(\mathbf{E})\) ）为纯的的必要充分条件是：对所有 \(u^{*} \in \bigwedge^{p-1}(\mathbf{E}^{*})\) ，

\[
(u ^ {*} \lrcorner z) \wedge z = 0.\tag{81}
\]

情形 \(p = 0\) 是平凡的，我们假设 \(p \geqslant 1\) 。若 \(z = x_1 \wedge \cdots \wedge x_p\) ，则公式 (68)（第 9 号）取 \(n = p - 1\) 表明 \(u^* \lrcorner z\) 是诸 \(x_i\) ( \(1 \leqslant i \leqslant p\) )的线性组合，由此得出 (81)。另一方面，若子空间 \(M_z\) （与 \(z\) 相伴）的维数 \(>p\) ，考虑此子空间的一组基 \((e_j)_{1 \leqslant j \leqslant n}\) （ \(n >p\) ）。由第 11 号命题 13 可知，每个 \(e_j\) 都具有 \(u^* \lrcorner z\) 的形式（对某个 \(u^* \in \bigwedge^{p-1}(E^*)\) ），因此关系式 (81) 蕴含 \(e_j \wedge z = 0\) （对 \(1 \leqslant j \leqslant n\) ）。由此可知，在表达式 \(z = \sum_{H} a_H e_H\) （其中 H 遍历 {[}1, n{]} 的含 \(p\) 个元素的子集的集合）中，所有系数 \(a_H\) 均为零，从而 \(z = 0\) ，与假设矛盾。

命题 16 的判据等价于当 \(u^*\) 遍历 \(\bigwedge^{p-1}(\mathbf{E}^*)\) 的一组基时写出条件 (81)。特别地，假设 \(\mathbf{E}\) 是有限维的，维数为 \(n\) ，并设 \((e_i)_{1\leq i\leq n}\) 是 \(\mathbf{E}\) 的一组基。则条件 (81) 等价于条件

\[
(82-(J, H)) \quad \langle e _ {\mathrm{J}} ^ {*}, (e _ {\mathrm{H}} ^ {*} \lrcorner z) \wedge z \rangle = 0
\]

对所有子集 J、H（ \([1, n]\) 的子集，满足 \(\text{Card}(J) = p + 1\) 和 \(\text{Card}(H) = p - 1\) ）成立。现在，若 I 和 \(I'\) 是 \([1, n]\) 的两个含 p 个元素的子集，则第 10 号的公式 (76) 以及 § 7 第 8 号的乘法表 (20) 表明

\[
\langle e _ {\mathrm{J}} ^ {*}, \langle e _ {\mathrm{H}} ^ {*} \lrcorner e _ {\mathrm{I}}) \wedge e _ {\mathrm{I} ^ {\prime}} \rangle = 0
\]

除非存在 \(i \in [1, n]\) ，使得 \(\mathbf{I} - \mathbf{H} = \{i\}\) 和 \(\mathbf{J} - \mathbf{I}' = \{i\}\) ，在这种情况下

\[
\langle e _ {\mathrm{J}} ^ {*}, (e _ {\mathrm{H}} ^ {*} \lrcorner e _ {\mathrm{I}}) \wedge e _ {\mathrm{I} ^ {\prime}} \rangle = (- 1) ^ {(p - 1) (p - 2) / 2} \varepsilon_ {i, \mathrm{J,H}}\tag{83}
\]

（其中 \(\varepsilon_{i,J,H} = \rho_{(i),H}\rho_{(i),I'}\) ；那么可以说，对 \(i\in J\cap \complement H\) , \(\varepsilon_{i,J,H}\) 等于 \(+1\) ，如果 \(J\) 中 \(< i\) 的元素个数与 \(H\) 中 \(< i\) 的元素个数有相同的奇偶性；否则等于 \(-1\) 。）

由此立即得出，若记 \(z = \sum_{\mathbf{I}} a_{\mathbf{I}} e_{\mathbf{I}}\) ，其中 \(\mathbf{I}\) 遍历 \([1, n]\) 的含 p 个元素的子集的集合，则关系 \((82-(J, H))\) 等价于关系

\[
(84 - (\mathrm{J}, \mathrm{H})) \quad \sum_ {i \in J \cup \complement_ {\mathrm{H}}} \varepsilon_ {i, J, \mathrm{H}} a _ {J - \{i \}} a _ {\mathrm{H} \cup \{i \}} = 0.
\]

关系 (84) 称为格拉斯曼关系；因此（当 J 遍历含 \(p + 1\) 个元素的子集的集合而 H 遍历含 \(p - 1\) 个元素的子集的集合（ \([1, n]\) 的）时），这些关系是元素 \(z \neq 0\) （属于 \(\bigwedge^{p}(\mathbf{E})\) ）为纯的必要充分条件。

注意，关系 (84) 并非独立的。例如，对 \(n = 4\) 和 \(p = 2\) ，格拉斯曼关系化简为单一关系

\[
a _ {1 2} a _ {3 4} - a _ {1 3} a _ {2 4} + a _ {1 4} a _ {2 3} = 0.\tag{85}
\]

设 \(D_p(E)\) 是 \(\bigwedge^p(E)\) 中由纯 \(p\) -向量组成的子集；显然 \(D_p(E)\) 关于 \(u\) 和 \(v\) 之间的等价关系：“存在 \(\lambda \in K^*\) 使得 \(v = \lambda u\)”是饱和的，且两个元素 \(u, v\) （ \(D_p(E)\) 的）在此关系下等价，当且仅当与它们相伴的子空间 \(M_u\) 和 \(M_v\) （ \(E\) 的）相同。因此我们由此得到一个从由 \(p\) 维向量子空间（ \(E\) 的）组成的集合到像 \(G_p(E)\) （即 \(D_p(E)\) 在射影空间 \(P(\bigwedge^p(E))\) （与 \(\bigwedge^p(E)\) 相伴）中的像）上的典范双射。子集 \(G_p(E)\) （ \(P(\bigwedge^p(E))\) 的）称为指标为 \(p\) 的向量空间 \(E\) 的格拉斯曼簇。当 \(E\) 有限维且 \((e_i)_{1 \leq i \leq n}\) 是 \(E\) 的一个基时，指标为 \(p\) 的格拉斯曼簇是 \(P(\bigwedge^p(E))\) 中使一个齐次坐标系 \((a_I)\) （相对于基 \((e_I)\) （ \(\bigwedge^p(E)\) 的基））满足格拉斯曼关系 (84) 的点的集合。

当 \(\mathbf{E} = \mathbf{K}^n\) 时，我们有时写 \(\mathbf{G}_{n,p}(\mathbf{K})\) 以代替 \(\mathbf{G}_p(\mathbf{K}^n)\) ，于是 \(\mathbf{G}_{n,1}(\mathbf{K}) = \mathbf{P}_{n-1}(\mathbf{K})\) 。映射 \(\mathbf{M} \mapsto \mathbf{M}^0\) 把每个 \(p\) 维子空间（ \(\mathbf{K}^n\) 的）对应到 \(\mathbf{E}^*\) 中的正交子空间（它借助于与 \(\mathbf{K}^n\) 的典范基对偶的基的选择而与 \(\mathbf{K}^n\) 等同），因而定义了 \(\mathbf{G}_{n,p}(\mathbf{K})\) 到 \(\mathbf{G}_{n,n-p}(\mathbf{K})\) 上的一个典范双射；命题 15 表明，这个双射是 \(\mathbf{G}_{n,p}(\mathbf{K})\) 上的这样一个限制：射影空间 \(\mathbf{P}(\bigwedge^p (\mathbf{K}^n))\) 到射影空间 \(\mathbf{P}(\bigwedge^{n-p}(\mathbf{K}^n))\) 的一个典范同构的限制。

\BourbakiAppendixSection{交错代数。八元数}

\subsection{交错代数}

设 A 是一个交换环，F 是一个（未必结合的）A-代数。对于 F 中任意三个元素 x, y, z，我们记

\[
a (x, y, z) = x (y z) - (x y) z\tag{1}
\]

（ \(x, y, z\) 的结合子）； \(a\) 显然是 \(\mathbf{F} \times \mathbf{F} \times \mathbf{F}\) 到 \(\mathbf{F}\) 中的一个 A-三线性映射。

引理 1. 对所有 \(p, q, r, s\) （属于代数 \(\mathbf{F}\) ），

\[
a (p q, r, s) - a (p, q r, s) + a (p, q, r s) = a (p, q, r) s + p a (q, r, s).\tag{2}
\]

验证由定义 (1) 立即得出。

命题 1. 对 A-代数 F，以下条件等价：

\begin{enumerate}
\def\labelenumi{(\alph{enumi})}
\item
  对 \(x, y\) （ \(\mathbf{F}\) 中元素的每个有序对），由 \(x\) 和 \(y\) 生成的子代数是结合的。
\item
  三线性映射 \((x, y, z) \mapsto a(x, y, z)\) 是交错的（§ 7，第 3 号）。
\item
  对 \(x, y\) （ \(\mathbf{F}\) 中元素的每个有序对）， \(x^2y = x(xy)\) 和 \(yx^2 = (yx)x\) 。
\end{enumerate}

显然（a）蕴含（c）。我们证明（c）蕴含（b）：根据定义（§ 7，第 3 号），要证明（b），只需验证 \(a(x, x, y) = 0\) 和 \(a(x, y, y) = 0\) ，而这正是（c）。

为了证明(b)蕴含(a)，我们使用以下4个引理：

引理 2. 设 E 是一个 A-代数，使得三线性映射 \((x,y,z)\mapsto a(x,y,z)\) 是交错的，S 是 E 的一个生成系，U 是 E 的一个包含 S 的子 A-模，并且满足 \(sU\subset U\) 和 \(Us\subset U\) （对所有 \(s\in S\) ）。则 U = E。

集合 \(U'\) 由 \(x \in E\) （满足 \(xU \subset U\) 和 \(Ux \subset U\) ）组成，它显然是 E 的一个子 A-模，且根据假设它包含 S。另一方面，对 \(U'\) 中的 x, y 和 \(u \in U\) ，根据假设

\[
(x y) u = x (y u) + a (x, y, u) = x (y u) - a (x, u, y) = x (y u) - (x u) y + x (u y) \in \mathrm{U};
\]

过渡到反代数，我们类似地有 \(u(xy) \in \mathbf{U}\) 。因此 \(\mathbf{U}'\) 是 \(\mathbf{E}\) 的一个子代数，并且由于它包含 \(\mathbf{S}\) , \(\mathbf{U}' = \mathbf{E}\) 。因此 \(\mathbf{EU} \subset \mathbf{U}\) ，更有 \(\mathbf{UU} \subset \mathbf{U}\) ，这证明了 \(\mathbf{U}\) 是 \(\mathbf{E}\) 的一个子代数；由于它包含 \(\mathbf{S}\) , \(\mathbf{U} = \mathbf{E}\) ，这就证明了引理。

F 的子集 H 称为强结合的，如果当 u、v、w 中至少有两个元素属于 H 时 \(a(u, v, w) = 0\) 。

引理 3. 假设映射 \(a\) 是交错的。如果 \(\mathbf{H}\) 是 \(\mathbf{F}\) 的一个强结合子集，则 \(\mathbf{F}\) 的由 \(\mathbf{H}\) 生成的子代数是强结合的。

由于 F 的强结合子集的集合是归纳的，只需证明：若 H 是 F 的一个极大强结合子集，则 H 是 F 的一个子代数。由于 H 显然是 F 的一个子 A-模，只需验证对 H 中任意两个元素 u、v, \(H \cup \{uv\}\) 也是强结合的，因为根据 H 的定义，这将蕴含 \(uv \in H\) 。现在，对所有 \(z \in H\) 以及所有 \(t \in F\) ，由 (2)

\[
a (u v, t, z) - a (u, v t, z) + a (u, v, t z) = 0
\]

因为 \(\mathbf{H}\) 是强结合的；由于 \(u, v, z\) 属于 \(\mathbf{H}\) ，同时

\[
a (u, v t, z) = a (u, v, t z) = 0,
\]

，因此 \(a(uv, t, z) = 0\) 。利用 a 是交错的这一事实，这表明 \(a(p, q, r) = 0\) （只要 p、q、r 这三个元素中至少有两个属于 \(H \cup \{uv\}\) ），由此即得引理。

引理 4. 假设映射 \(a\) 是交错的。那么，对所有 \(x \in \mathbf{F}\) ， \(\mathbf{F}\) 的由 \(x\) 生成的子代数是强结合的。

\(a(u,v,w) = 0\) （只要 \(u, v, w\) 这三个元素中有两个等于 \(x\) ），于是只需应用引理 3。

引理 5. 假设映射 \(a\) 是交错的，并设 \(\mathbf{X}, \mathbf{Y}\) 是 \(\mathbf{F}\) 的两个强结合子代数。则 \(\mathbf{E}\) 的由 \(\mathbf{X} \cup \mathbf{Y}\) 生成的子代数是结合的。

设 Z 为 \(z \in E\) （满足 \(a(u, v, z) = 0\) ，对所有 \(u \in X\) 和 \(v \in Y\) ）的集合，它显然是一个包含 X 和 Y 的子 A-模，因为 X 和 Y 是强结合的；由引理 2，只需验证：对 \(u \in X\) 和 \(v \in Y\) , \(uZ \subset Z\) , \(vZ \subset Z\) , \(Zu \subset Z\) 和 \(Zv \subset Z\) 。现在，对 X 中的 \(u, u'\) ， \(v \in Y\) 和 \(z \in Z\) ，由 (2)

\[
a (u ^ {\prime} u, z, v) - a (u ^ {\prime}, u z, v) + a (u ^ {\prime}, u, z v) = a (u ^ {\prime}, u, z) v + u ^ {\prime} a (u, z, v) = 0
\]

这是由于 X 是强结合的以及 Z 的定义。但鉴于 X 是强结合的， \(a(u', u, zv) = 0\) ，并且由于 \(u'u \in X\) ，按 Z 的定义有 \(a(u'u, z, v) = 0\) 。因此 \(a(u', uz, v) = 0\) ，这表明 \(uZ \subset Z\) 。现在将 (2) 应用于 \((p, q, r, s) = (v, z, u, u')\) ，我们类似地得到 \(Zu \subset Z\) 。交换 X 与 Y 的角色，并利用 a 是交错的事实，我们得到 \(vZ \subset Z\) 和 \(Zv \subset Z\) ；由此得引理。

现在，为证明命题 1 中 (b) 蕴含 (a)，只需取 \(\mathbf{X} = \{x\}\) 和 \(\mathbf{Y} = \{y\}\) ，并利用引理 4。

定义 1. 代数 \(\mathbf{F}\) 称为交错的，如果它满足命题 1 的等价条件。

结合代数显然是交错代数。在第 3 号中，我们将给出一个交错但不结合的代数的例子。

若 \(\mathbf{F}\) 是一个交错 A-代数，则每个 A-代数 \(\mathbf{F} \otimes_{\mathbf{A}} \mathbf{A}'\) （由 \(\mathbf{F}\) 通过标量扩张（§ 1，第 5 号）得到）都是一个交错 \(\mathbf{A}'\) -代数，这由命题 1 的条件 (b) 得出。

\subsection{交错凯莱代数}

命题 2. 设 A 为一个环，F 为一个凯莱 A-代数，e 为其单位元， \(s: x \mapsto \bar{x}\) 为其共轭，N: F \(\to\) A 为其凯莱范数（§ 3，第 4 号）。

\begin{enumerate}
\def\labelenumi{(\roman{enumi})}
\item
  \(\mathbf{F}\) 为交错代数的必要充分条件是：对 \(x, y\) （ \(\mathbf{F}\) 中元素的每个有序对），有 \(x^2y = x(xy)\) 。
\item
  若 \(\mathbf{F}\) 是交错的，则 \(\mathbf{N}(xy) = \mathbf{N}(x)\mathbf{N}(y)\) 对所有 \(x,y\) （属于 \(\mathbf{F}\) ）成立。
\item
  假设 \(\mathbf{F}\) 是交错的。元素 \(x \in \mathbf{F}\) 可逆的必要充分条件是 \(\mathbf{N}(x)\) 在 \(\mathbf{Ae}\) 中可逆；此时 \(x\) 的逆元是唯一的，且等于 \(\mathbf{N}(x)^{-1}\bar{x}\) ；把它记作 \(x^{-1}\) ，
\end{enumerate}

\[
x ^ {- 1} (x y) = x \left(x ^ {- 1} y\right) = y
\]

对所有 \(y \in \mathbf{F}\) 成立。

条件 \(x^{2}y = x(xy)\) 对于 F 是交错代数显然是必要的（第 1 号，命题 1）。反之，若它对 F 的元素的每个有序对都成立，将其应用于 \(\bar{x}\) 和 \(\bar{y}\) ，就得到 \(\bar{x}^{2}\bar{y} = \bar{x}(\bar{x}\bar{y})\) ；对这一关系式应用共轭 s，我们得到 \(yx^{2} = (yx)x\) ，因此第 1 号命题 1 的条件 (c) 得以满足。

显然 \(a(e, x, y) = 0\) 对所有 \(x, y\) （属于 \(\mathbf{F}\) ）成立。若 \(\mathbf{F}\) 是交错的，则由命题 1（第 1 号）推出，子代数 \(\mathbf{G}\) （ \(\mathbf{F}\) 中由 \(e, x\) 和 \(y\) 生成的子代数）是结合的。由于 \(\bar{x} = -x + \mathrm{T}(x) \in -x + \mathrm{Ae}, \bar{x} \in \mathbf{G}\) ，类似地 \(\bar{y} \in \mathbf{G}\) 。于是 \(\mathrm{N}(xy) = (xy)(\overline{xy}) = xy.\bar{y}.\bar{x} = \mathrm{N}(y)x\bar{x} = \mathrm{N}(y)\mathrm{N}(x)\) ，这里利用了 \(\mathrm{N}(y) \in \mathrm{Ae}\) 。这就证明了 (ii)。

最后我们证明 (iii)。若 \(\mathrm{N}(x)\) 在 Ae 中可逆，且记 \(x' = \mathrm{N}(x)^{-1}\bar{x}\) ，则 \(xx' = x'x = e\) ，因为 \(\mathrm{N}(x) = x\bar{x} = \bar{x}x\) 。反之，若 x 有左逆 \(x''\) ，则 \(\mathrm{N}(x'')\mathrm{N}(x) = \mathrm{N}(e) = e\) （由 (ii)），故 \(\mathrm{N}(x)\) 在 Ae 中可逆；此外，由于 \(x' = \mathrm{N}(x)^{-1}\bar{x}\) 属于由 x 和 e 生成的子代数，元素 \(x, x', x''\) 都属于由 x、 \(x''\) 和 e 生成的结合子代数，因此 \(x'' = x''(xx') = (x''x)x' = x'\) ，唯一性由此得证。公式 \(x^{-1}(xy) = x(x^{-1}y) = y\) 由如下事实得出： \(x^{-1}\) 、x 和 y 是由 x、y 和 e 生成的子代数的元素，而该子代数是结合的。

命题 3. 设 E 是一个凯莱 A-代数， \(\gamma\) 是 A 的一个元素，F 是 E 的由 \(\gamma\) 与 E 的共轭所定义的凯莱扩张（§ 2，第 5 号，命题 5）。F 为交错代数的必要充分条件是 E 为结合代数。

设 \(u = (x, y)\) , \(v = (x', y')\) 是 \(\mathbf{F}\) 的两个元素（其中 \(x, y, x', y'\) 属于 \(\mathbf{E}\) ）。则（§ 2，第 5 号，公式 (27)）

\[
\left\{ \begin{array}{l} u ^ {2} v = ((x ^ {2} + \gamma \bar {y} y) x ^ {\prime} + \gamma \bar {y} ^ {\prime} (y \bar {x} + y x), (y \bar {x} + y x) \bar {x} ^ {\prime} + y ^ {\prime} (x ^ {2} + \gamma \bar {y} y)) \\ u (u v) = (x (x x ^ {\prime} + \gamma \bar {y} ^ {\prime} y) + \gamma (x ^ {\prime} \bar {y} + \bar {x}. \bar {y} ^ {\prime}) y, y (\bar {x} ^ {\prime} \bar {x} + \gamma \bar {y} y ^ {\prime}) + (y \bar {x} ^ {\prime} + y ^ {\prime} x) x). \end{array} \right.\tag{3}
\]

利用 \(\bar{y}y\) 和 \(\bar{x} + x\) 属于 Ae 这一事实，考察这些公式表明，E 的结合性蕴含 \(u^{2}v = u(uv)\) ，从而 F 是交错的（命题 2）。反之，若 F 是交错的，则方程 \(u^{2}v = u(uv)\) 在 \(y' = 0\) 时给出

\[
(y \bar {x} + y x) \bar {x} ^ {\prime} = y (\bar {x} ^ {\prime} \bar {x}) + (y \bar {x} ^ {\prime}) x.
\]

现在左边等于 \((y\mathrm{T}(x))\bar{x}' = y(\bar{x}'\mathrm{T}(x)) = y(\bar{x}'x + \bar{x}'\bar{x})\) ；

与右边比较，我们得到 \((y\bar{x}^{\prime})x = y(\bar{x}^{\prime}x)\) ，这证明了 E 的结合性，因为 x、y 和 \(\bar{x}^{\prime}\) 是 E 的任意元素。

\subsection{八元数}

设 E 是 A 上类型为 \((\alpha, \beta, \gamma)\) 的四元数代数（§ 2，第 5 号，例 2），并设 \(\delta \in A\) 。E 的由 \(\delta\) 与 E 的共轭所定义的凯莱扩张 F 称为 A 上的八元数代数，并称其类型为 \((\alpha, \beta, \gamma, \delta)\) 。由第 2 号的命题 3，F 是交错代数。它有一个由 8 个元素组成的基 \((e_{i})_{0 \leq i \leq 7}\) ，定义为

\[
\begin{array}{l l l l} e _ {0} = (e, 0), & e _ {1} = (i, 0), & e _ {2} = (j, 0), & e _ {3} = (k, 0) \\ e _ {4} = (0, e), & e _ {5} = (0, i), & e _ {6} = (0, j), & e _ {7} = (0, k) \end{array}
\]

（其中 \((e, i, j, k)\) 是 E 的基，定义于 loc. cit.；显然 \(e_0\) （也记作 \(e\) ）是 F 的单位元。若 \(u = \sum_{i=0}^{7} \xi_i e_i\) 是 F 的一个元素（其中 \(\xi_i \in A\) ），则 § 2 第 5 号的公式 (23)、(24) 和 (31) 给出八元数 \(u\) 的共轭、迹和范数

\[
\left\{ \begin{array}{c} \bar {u} = (\xi_ {0} + \beta \xi_ {1}) e _ {0} - \sum_ {i = 1} ^ {7} \xi_ {i} e _ {i} \\ \mathrm{T} _ {\mathrm{F}} (u) = 2 \xi_ {0} + \beta \xi_ {1} \\ \mathrm{N} _ {\mathrm{F}} (u) = \xi_ {0} ^ {2} + \beta \xi_ {0} \xi_ {1} - \alpha \xi_ {1} ^ {2} - \gamma (\xi_ {2} ^ {2} + \beta \xi_ {2} \xi_ {3} - \alpha \xi_ {3} ^ {2}) \\ - \delta (\xi_ {4} ^ {2} + \beta \xi_ {4} \xi_ {5} - \alpha \xi_ {5} ^ {2}) + \gamma \delta (\xi_ {6} ^ {2} + \beta \xi_ {6} \xi_ {7} - \alpha \xi_ {7} ^ {2}). \end{array} \right.\tag{4}
\]

现在设 \(u = (x, y)\) , \(u' = (x', y')\) 和 \(u'' = (x'', y'')\) 是三个八元数（其中元素 \(x, x', x'', y, y', y''\) 属于 E）。§ 2 第 5 号的公式 (24) 和 (27) 给出

\[
\begin{array}{r l} & {\mathrm{T} _ {\mathrm{F}} ((u u ^ {\prime}) u ^ {\prime \prime}) = \mathrm{T} (x x ^ {\prime} x ^ {\prime \prime}) + \delta \mathrm{T} (\bar {y} ^ {\prime} y x ^ {\prime \prime}) + \delta \mathrm{T} (\bar {y} ^ {\prime \prime} y \bar {x} ^ {\prime}) + \delta \mathrm{T} (\bar {y} ^ {\prime \prime} y ^ {\prime} x)} \\ & {\mathrm{T} _ {\mathrm{F}} (u (u ^ {\prime} u ^ {\prime \prime})) = \mathrm{T} (x x ^ {\prime} x ^ {\prime \prime}) + \delta \mathrm{T} (x ^ {\prime \prime} \bar {y} ^ {\prime} y) + \delta \mathrm{T} (\bar {x} ^ {\prime} \bar {y} ^ {\prime \prime} y) + \delta \mathrm{T} (x \bar {y} ^ {\prime \prime} y ^ {\prime})} \end{array}
\]

（其中 T 表示 E 中的迹，并利用了 E 是结合的这一事实）。由于 \(\mathrm{T}(xy) = \mathrm{T}(yx)\) 对所有四元数 x、y 成立（§ 2，第 4 号，公式 (17)），由此可得

\[
\mathrm{T} _ {\mathrm{F}} ((u u ^ {\prime}) u ^ {\prime \prime}) = \mathrm{T} _ {\mathrm{F}} (u (u ^ {\prime} u ^ {\prime \prime})).\tag{5}
\]

我们特别研究类型为 \((-1, 0, -1, -1)\) 的八元数；此时公式（4）简化为

\[
\left\{ \begin{array}{c} {\bar {u}} = \xi_ {0} e _ {0} - \sum_ {i = 1} ^ {7} \xi_ {i} e _ {i} \\ \mathrm{T} _ {\mathrm{F}} (u) = 2 \xi_ {0} \\ \mathrm{N} _ {\mathrm{F}} (u) = \sum_ {i = 0} ^ {7} \xi_ {i} ^ {2}. \end{array} \right.\tag{6}
\]

*如果我们取 A 为实数域 R，则 R 上类型为 \((-1,0,-1,-1)\) 的八元数称为凯莱八元数（或 octaves）。由第 2 号的命题 2 (ii) 可知，每个 \(\neq0\) 的凯莱八元数都是可逆的。*

命题 4. 设 F 是 A 上类型为 \((-1,0,-1,-1)\) 的一个八元数代数。存在一个以二元素域 Z/2Z 为标量域的三维向量空间 V，以及一个双射 \(\lambda\mapsto e_{\lambda}^{\prime}\) （从 V 到基 \((e_{i})_{0\leq i\leq7}\) 上的双射），使得

\[
e _ {0} ^ {\prime} = e _ {0}, \quad e _ {\lambda} ^ {\prime} e _ {\mu} ^ {\prime} = \pm e _ {\lambda + \mu} ^ {\prime}\tag{7}
\]

对所有 \(\lambda, \mu\) （属于 V）。要使 \(e_{\lambda}'(e_{\mu}'e_{\nu}') = (e_{\lambda}'e_{\mu}')e_{\nu}'\) ，只需 \(\lambda, \mu, \nu\) （ V 中的）在 Z/2Z 上线性相关即可；若 \(2 \neq 0\) （在 A 中），则此条件是必要的。

我们保留本节开头的记号。由 § 2 第 5 号的公式 (33) 可知，由元素 \(\pm e_0\) , \(\pm e_1\) , \(\pm e_2\) , \(\pm e_3\) 组成的集合 S 在乘法下是稳定的。此外，对 \(x, y, y'\) （属于 E），由 § 2 第 5 号的公式 (22)，

\[
(x, 0) (0, y ^ {\prime}) = (0, y ^ {\prime} x), \quad (0, y ^ {\prime}) (x, 0) = (0, y ^ {\prime} \bar {x}), \quad (0, y) (0, y ^ {\prime}) = (- \bar {y} ^ {\prime} y, 0)\tag{8}
\]

使得由元素 \(\pm e_{i}\) ( \(0 \leqslant i \leqslant 7\) ) 组成的集合 T 在乘法下是稳定的；此外，其乘法表与环 A 无关。

特别地，设 \(\mathbf{A}''\) 为含两个元素的域 \(\mathbf{Z}/2\mathbf{Z}\) ，设 \(\mathbf{E}''\) 为 \(\mathbf{A}''\) 上类型为 (1, 0, 1) 的四元数代数， \(\mathbf{F}''\) 为 \(\mathbf{A}''\) 上类型为 (1, 0, 1, 1) 的八元数代数；令 \((e_i'')_{0 \leq i \leq 7}\) 为 \(\mathbf{F}''\) 的按上述方式构成的基。由于 \(-e_i'' = e_i''\) ，集合 \(\mathbf{T}''\) （由 \(e_i''\) 组成）有 8 个元素且在乘法下封闭；此外，由上述内容立即得出，映射 \(\theta: \mathbf{T} \to \mathbf{T}''\) （满足 \(\theta(e_i) = \theta(-e_i) = e_i''\) ，对所有 \(0 \leq i \leq 7\) ）是乘法同态。此外，四元数代数 \(\mathbf{E}''\) 在此情形下是交换的，因而 \(\mathbf{F}''\) 是结合的（ \(\S 2\) ，第 5 号，命题 5）；进一步， \(\mathbf{F}''\) 中的共轭在此情形下是恒等映射。因此 \(\mathbf{T}''\) 是一个群，且公式 (8) 表明它是交换的；这些公式以及 \(\S 2\) 第 5 号的公式 (33) 表明 \(\mathbf{T}''\) 的每个元素的平方都是单位元。若用 V 表示写成加法群后的群 \(\mathbf{T}''\) ，则可在 V 上赋予唯一的向量空间结构，其标量域为 \(\mathbf{Z}/2\mathbf{Z}\) ，其维数必然为 3，因为

\[
\operatorname{Card} (\mathbf {V}) = 8 = 2 ^ {\dim (\mathbf {V})}.
\]

对所有 \(\lambda \in V\) ，随后设 \(e_{\lambda}'\) 表示 \((e_i)_{0 \leq i \leq 7}\) 中满足 \(\theta(e_{\lambda}') = \dot{\lambda}\) 的元素；于是 \(e_0' = e_0\) ；此外，由于 \(\theta\) 是同态，且关系 \(\theta(x) = \theta(y)\) 等价于 \(x = \pm y\) ，故有 \(e_{\lambda}' e_{\mu}' = \pm e_{\lambda + \mu}'\) 。若 \(\lambda, \mu, \nu\) 在 \(\mathbf{Z}/2\mathbf{Z}\) 上线性无关，它们就构成 \(V\) 的一组基，从而所有元素 \(e_i\) ( \(0 \leqslant i \leqslant 7\) ) 都将属于由 \(e_{\lambda}', e_{\mu}'\) 和 \(e_{\nu}'\) 生成的子代数；当 \(2 \neq 0\) （在 \(A\) 中）时， \(e_{\lambda}'(e_{\mu}' e_{\nu}') = (e_{\lambda}' e_{\mu}') e_{\nu}'\) 是不可能的，因为那样 \(F\) 就是结合的，从而 \(E\) 是交换的（ \(\S 2\) ，第 5 号，命题 5），而这与 \(\S 2\) 第 5 号的关系式 (33) 相矛盾。另一方面，如果 \(\lambda, \mu, \nu\) 在 V 中线性

相关，则这三个元素 \(e_{\lambda}^{\prime}\) , \(e_{\mu}^{\prime}\) , \(e_{\nu}^{\prime}\) 属于 F 的一个由 2 个生成元生成的子代数，因而该子代数是结合的（第 1 号，命题 1）；结论由此得出。

注记。由于 \(\bar{e}_{\lambda}^{\prime} = -e_{\lambda}^{\prime}\) （对 \(\lambda \neq 0\) ），

\[
e _ {\lambda} ^ {\prime 2} = - e _ {0} \quad \text { for } \lambda \neq 0,
\]

\[
e _ {\mu} ^ {\prime} e _ {\lambda} ^ {\prime} = - e _ {\lambda} ^ {\prime} e _ {\mu} ^ {\prime} \quad \text { for } \lambda \neq 0, \mu \neq 0 \text { and } \mu \neq \lambda .
\]

\BourbakiExercisesHeading

\BourbakiExercisesSection

\begin{enumerate}
\def\labelenumi{\arabic{enumi}.}
\tightlist
\item
  设 E 是（交换）域 K 上的一个代数，且 L 是 K 的一个交换扩域。证明：若代数 \(\mathrm{E}_{(L)}\) 具有单位元，则 E 亦然（参见 II，§ 3，第 5 号，命题 6）。
\end{enumerate}

\BourbakiExercisesSection

\begin{enumerate}
\def\labelenumi{\arabic{enumi}.}
\tightlist
\item
  设 E 是一个 A-代数，具有由两个元素 \(e_1, e_2\) 组成的基和一个单位元 \(e = \lambda e_1 + \mu e_2\) 。
\end{enumerate}

\begin{enumerate}
\def\labelenumi{(\alph{enumi})}
\item
  证明理想 \(a\) （ A 的由 \(\lambda\) 和 \(\mu\) 生成的理想）就是整个环 A（在相反的情形，注意 \(\mathrm{E} / \mathrm{aE} = \mathrm{E} \otimes_{\mathbf{A}} (\mathrm{A} / \mathrm{a}) = \{0\}\) 将与如下事实矛盾： \(\mathrm{A} / \mathrm{a} \neq \{0\}\) 且 E 是自由 A-模）。
\item
  若 \(\alpha, \beta\) 是 A 的两个元素，使得 \(\alpha\lambda + \beta\mu = 1\) ，证明元素 \(e\) 和 \(u = \beta e_1 - \alpha e_2\) 也构成 E 的一组基，并由此推出 E 是一个二次 A-代数。
\end{enumerate}

\begin{enumerate}
\def\labelenumi{\arabic{enumi}.}
\setcounter{enumi}{1}
\item
  证明 \(\mathbf{Z}\) 上的每个二次代数同构于类型为 \((0, n)\) 的代数（其中 \(n \in \mathbf{N}\) ），或类型为 \((m, 0)\) 的代数（其中 \(m\) 非平方数），或类型为 \((m, 1)\) 的代数（其中 \(m\) 是不具有形式 \(k(k - 1)\) 的整数）；此外，这些代数中任意两个都不同构。
\item
  \begin{enumerate}
  \def\labelenumii{(\alph{enumii})}
  \tightlist
  \item
    设 K 是一个交换域。证明：K 为 K 上类型为 \((\alpha, \beta, \gamma)\) 的四元数代数的中心的必要充分条件是以下情形之一成立：(1) \(\beta\gamma \neq 0\) ；(2) \(\gamma = 0\) , \(\beta \neq 0\) 且 \(\beta^{2} + 4\alpha \neq 0\) ；(3) \(\beta = 0\) , \(\alpha \neq 0\) 或 \(\gamma \neq 0\) ，且 \(2 \neq 0\) 在 K 中。
  \end{enumerate}
\end{enumerate}

\begin{enumerate}
\def\labelenumi{(\alph{enumi})}
\setcounter{enumi}{1}
\item
  设 K 是一个交换域，使得 \(2 \neq 0\) 在 K 中；证明 K 上类型为 \((1, \gamma)\) 的四元数代数同构于矩阵代数 \(\mathbf{M}_{2}(\mathbf{K})\) （考虑该代数由元素 \(\frac{1}{2}(1 + i)\) , \(\frac{1}{2}(1 - i)\) , \((j + k/2\gamma)\) , \(\frac{1}{2}(j - k)\) 组成的基）。
\item
  设 K 是一个交换域，使得 \(2 \neq 0\) 在 K 中，并设 E 为 K 上类型为 \((\alpha, \gamma)\) 的四元数代数。四元数 \(z \in E\) 称为纯的，如果 \(\overline{z} = -z\) （或等价地 \(T(z) = 0\) ）。证明：若 z 是纯四元数，则每个四元数 \(tzt^{-1}\) （其中 \(t \in E\) 可逆）也是纯的。
\item
  若 \(u, v\) 是任意两个四元数，则 \(uv - vu\) 是纯的。由此推出，若 \(u, v, w\) 是任意三个四元数，则
\end{enumerate}

\[
(u v - v u) ^ {2} w = w (u v - v u) ^ {2}.
\]

\begin{enumerate}
\def\labelenumi{(\alph{enumi})}
\setcounter{enumi}{4}
\tightlist
\item
  证明：若 E 中存在四元数 \(z \neq 0\) 使得 \(\mathrm{N}(z) = 0\) ，则也存在纯四元数 \(z' \neq 0\) 使得 \(\mathrm{N}(z') = 0\) 。（注意，在公式 (33)（第 5 号）的记号下，若 \(\alpha\) 不是 K 中的平方数，则 E 中存在 \(z \neq 0\) 使得 \(\mathrm{N}(z) = 0\) 这一事实等价于存在元素 \(y \in \mathrm{K} + \mathrm{Ki}\) 使得 \(\gamma = \mathrm{N}(y)\) 。）
\end{enumerate}

\begin{enumerate}
\def\labelenumi{\arabic{enumi}.}
\setcounter{enumi}{4}
\tightlist
\item
  在满足 2 = 0 的交换域 K 上，每个类型为 \((\alpha, 0, \gamma)\) 的四元数代数 E 都是交换的，且每个 \(x \in E\) 的平方都属于 K；因此，子代数 \(\mathrm{K}(x)\) （由 E 的元素 \(x \in K\) 生成）是 K 上类型为 \((\lambda, 0)\) 的二次代数，且 E 是 \(\mathrm{K}(x)\) 上的二次代数。
\end{enumerate}

证明集合 C （由 \(x \in E\) （其平方等于 K 中某个元素的平方）组成）是 E 的一个向量子空间，其维数等于 1、2 或 4。若 C 的维数为 1（此时 C = K），则 E 是一个域。若 C 的维数为 2，则存在一个包含于 E 中的二次代数 \(\mathrm{K}(x)\) ，它是一个域，并且 E 在 \(\mathrm{K}(x)\) 上有一组由两个元素 1 和 u 组成的基，满足 \(u^{2} = 0\) ； \(y \in E\) （满足 \(y^{2} = 0\) ）的集合是 K 上的一个 2 维理想 a，且 E/a 同构于 \(\mathrm{K}(x)\) 。最后，若 C 的维数为 4（从而等于 E），则 \(y \in E\) （满足 \(y^{2} = 0\) ）的集合 a 是一个 3 维理想，且 E/a 同构于 K；E 中存在一个基 \((1, e_{1}, e_{2}, e_{3})\) ，满足 \(e_{1}^{2} = e_{2}^{2} = e_{3}^{2} = 0\) , \(e_{1}e_{2} = e_{3}\) , \(e_{1}e_{3} = e_{2}e_{3} = 0\) ； \(Ke_{3} = b\) 是 E 中唯一的 1 维理想，b 是 a 的零化子，且 a/b 是 E/b 的两个 1 维理想的直和，这两个理想相互零化。

\begin{enumerate}
\def\labelenumi{\arabic{enumi}.}
\setcounter{enumi}{5}
\tightlist
\item
  设 \(\mathbf{K}\) 是一个满足 \(2 \neq 0\) 的交换域， \(\mathbf{E}\) 是 \(\mathbf{K}\) 上类型为 \((0, 0, \gamma)\) 的四元数代数。
\end{enumerate}

\begin{enumerate}
\def\labelenumi{(\alph{enumi})}
\item
  若 \(\gamma\) 不是 \(\mathbf{K}\) 中的平方数，则 \(\mathbf{E}\) 中不存在 \(\mathbf{K}\) 上维数为 1 的左（相应地，右）理想；集合 \(a\) （由 \(x \in \mathbf{E}\) （满足 \(x^2 = 0\) ）组成）是 \(\mathbf{K}\) 上维数为 2 的双边理想，且 \(\mathbf{E} / a\) 是一个域，即 \(\mathbf{K}\) 上的一个二次代数。
\item
  若 \(\gamma\) 是一个平方数 \(\neq 0\) （在 \(\mathbf{K}\) 中），则 \(\mathbf{E}\) 中存在一组基 \((e_i)_{1\leq i\leq 4}\) ，其乘法表如下：
\end{enumerate}

\(e_1\)

\(e_2\)

\(e_3\)

\(e_4\)

\(e_1\)

\(e_1\)

0

\(e_3\)

0

\(e_2\)

0

\(e_2\)

0

\(e_4\)

\(e_3\)

0

\(e_3\)

0

0

\(e_4\)

\(e_4\)

0

0

0

 \(x \in E\) （满足 \(x^{2} = 0\) ）的集合 a 是一个维数为 2 的双边理想，它是两个相互零化的双边理想 \(Ke_{3}\) 和 \(Ke_{4}\) 的直和；后者是 E 中仅有的维数为 1 的（左或右）理想。商代数 E/a 同构于两个与 K 同构的域的乘积。

\begin{enumerate}
\def\labelenumi{(\alph{enumi})}
\setcounter{enumi}{2}
\tightlist
\item
  若 \(\gamma = 0\) ，则集合 \(a\) （由 \(x \in E\) （满足 \(x^2 = 0\) ）组成）是一个维数为 3 的双边理想； \(Kk = b\) 是 \(E\) 中仅有的维数为 1 的（左或右）理想；它是一个双边理想，即 \(a\) 的左、右零化子。在商代数 \(E/b\) 中， \(a/b\) 是两个相互零化的 1 维双边理想的直和；最后， \(E/a\) 是一个与 \(K\) 同构的域。
\end{enumerate}

\begin{enumerate}
\def\labelenumi{\arabic{enumi}.}
\setcounter{enumi}{6}
\tightlist
\item
  设 K 是一个交换域，其中 \(2 \neq 0\) 且E为K上的代数，其基由四个元素1、i、j、k组成（1为单位元），乘法表如下：
\end{enumerate}

\[
i ^ {2} = j ^ {2} = k ^ {2} = 1, \quad i j = j i = k, \quad j k = k j = i, \quad k i = i k = j.
\]

证明 E 同构于四个与 K 同构的域的乘积（考虑由元素 \((1 + \varepsilon i)(1 + \varepsilon'j)\) （其中 \(\varepsilon\) 和 \(\varepsilon'\) 等于 1 或 -1）组成的 E 的基）。

代数 E 是（域 K 上的）两个 2 阶循环群的乘积群的代数。将此推广到 n 个 2 阶循环群的乘积群的代数（参见 § 4，第 1 号，例 2）。

\begin{enumerate}
\def\labelenumi{\arabic{enumi}.}
\setcounter{enumi}{7}
\item
  四元数群 \(\mathfrak{Q}\) （I，第 6 节，习题 4）同构于八个四元数 \(\pm 1\) , \(\pm i\) , \(\pm j\) , \(\pm k\) （在类型为 \((-1, -1)\) 的四元数代数（定义于满足 \(2 \neq 0\) 的域上）中）所成的乘法群。证明群 \(\mathfrak{Q}\) 在满足 \(2 \neq 0\) 的域 K 上的代数同构于四个与 K 同构的域之积以及 K 上类型为 \((+1, -1)\) 的四元数代数（若 \(c\) 是 \(\mathfrak{Q}\) 中对应于四元数 \(-1\) 的元素，则 \(\mathfrak{Q}\) 的元素可写成 \(e\) , \(i, j, k, c, ci, cj, ck\) ；考虑由元素 \(\frac{1}{2}(e + c)\) , \(\frac{1}{2}(e - c)\) , \(\frac{1}{2}(e + c)i\) , \(\frac{1}{2}(e - c)i\) , \(\frac{1}{2}(e + c)j\) , \(\frac{1}{2}(e - c)j\) , \(\frac{1}{2}(e + c)k\) , \(\frac{1}{2}(e - c)k\) 组成的 E 的基）。
\item
  \begin{enumerate}
  \def\labelenumii{(\alph{enumii})}
  \tightlist
  \item
    证明 8 阶二面体群 \(\mathbf{D}_4\) （I，第 6 节，习题 4）在满足 \(2 \neq 0\) 的域 K 上的代数 E 同构于四个与 K 同构的域的乘积以及 K 上的 2 阶矩阵代数。（若 \(a, b\) 是 I，第 6 节，习题 4 中引入的 \(\mathbf{D}_4\) 的两个生成元，则 \(\mathbf{D}_4\) 的元素可写成 \(a^i b^j\) ，其中 \(0 \leqslant i \leqslant 3\) , \(0 \leqslant j \leqslant 1\) ；考虑由四个元素 \(\frac{1}{2}(e + a^2)\) , \(\frac{1}{2}(e - a^2)\) , \(\frac{1}{2}(a + a^3)\) , \(\frac{1}{2}(a - a^3)\) 以及这四个元素右乘 \(b\) 所组成的 E 的基；利用习题 4。）由此推出，若 \(-1\) 是 K 中的平方数，则互不同构的群 \(\mathfrak{Q}\) 和 \(\mathbf{D}_4\) 在 K 上的代数是同构的。
  \end{enumerate}
\end{enumerate}

\begin{enumerate}
\def\labelenumi{(\alph{enumi})}
\setcounter{enumi}{1}
\tightlist
\item
  类似地证明：6 阶二面体群 \(D_{3}\) 在满足 \(2 \neq 0\) 和 \(3 \neq 0\) 的域 K 上的代数 F 同构于两个与 K 同构的域的乘积以及 K 上的 2 阶矩阵代数（这里考虑由元素 \(e + a + a^{2}\) , \(a + a^{2} - 2e\) , \(a - a^{2}\) 以及这三个元素右乘 b 所组成的 F 的基）。
\end{enumerate}

\begin{enumerate}
\def\labelenumi{\arabic{enumi}.}
\setcounter{enumi}{9}
\item
  设 G 是一个群，H 是 G 的一个正规子群。证明：若 a 是群代数 \(\mathbf{A}^{(\mathbf{G})}\) 的由元素 ts - s（其中 t 遍历 H，s 遍历 G）生成的双边理想，则群代数 \(\mathbf{A}^{(\mathbf{G}/\mathbf{H})}\) 同构于商代数 \(\mathbf{A}^{(\mathbf{G})}/\mathbf{a}\) 。
\item
  设 A 为交换环，S 为具有单位元 e 的交换幺半群。假设在 A 上给定一个以 S 为算子域的外合成律 \((s, x) \mapsto x^{s}\) ，使得对所有 \(s \in S\) ，映射 \(x \mapsto x^{s}\) 都是环 A 的自同构，且 \((x^{s})^{t} = x^{st}\) 对所有 \(x \in A\) 以及 S 中的 s 和 t 成立（这蕴含 e 是所考虑的外合成律的单位算子）。则在 A-模 \(\mathrm{A}^{(\mathrm{S})}\) 上定义了一个乘法的内部合成律，其典范基记为 \((b_{\mathrm{s}})\) ，由关系式
\end{enumerate}

\[
\left(\sum_ {s} x _ {s} b _ {s}\right) \left(\sum_ {s} y _ {s} b _ {s}\right) = \sum_ {s} \left(\sum_ {t u = s} a _ {t, u} x _ {t} y _ {u} ^ {s}\right) b _ {s}
\]

其中 \(a_{s,t}\) 属于 A。

证明这个合成律与 \(\mathbf{A}^{(\mathbf{S})}\) 上的加法一起在此集合上定义了一个环结构，只要 \(a_{s,t}\) 满足条件

\[
a _ {s, t} a _ {s t, u} = a _ {t u} ^ {s} a _ {s, t u}
\]

对于 S 中的所有 s、t、u。若还满足以下条件，则元素 \(b_{e}\) 是这个环的单位元

\[
a _ {e, s} = a _ {s, e} = 1
\]

对所有 \(s \in S\) 成立；此时 A 可等同于 \(A^{(S)}\) 的一个子环，并且若用 C 表示 A 的由在所有自同构 \(x \mapsto x^s\) 下不变的元素组成的子环，则 \(A^{(S)}\) 是一个 C-代数，称为环 A 与幺半群 S 关于因子系 \((a_{s,t})\) 的交叉积。若把因子系 \((a_{s,t})\) 换成 \(\left(\frac{c_s c_t^s}{c_{st}} a_{s,t}\right)\) （其中，对所有 \(s \in S\) , \(c_s\) 是 A 的可逆元素，且 \(c_e = 1\) ），则所得到的新交叉积同构于由因子系 \((a_{s,t})\) 所定义的交叉积。

特别地，若取 A 为环 C 上的二次代数，且 S 为 2 阶循环群，例如 \(\{e, s\}\) ，使得 \(x^s = \bar{x}\) 对所有 \(x \in \mathbf{A}\) 成立，证明 A 与 S 的每个交叉积都是 C 上的四元数代数。

\begin{enumerate}
\def\labelenumi{\arabic{enumi}.}
\setcounter{enumi}{11}
\tightlist
\item
  设 I 为非空集合。证明自由广群 M(I) 的元素可以与序列 \((a_{i})_{1 \leqslant i \leqslant n}\) （其中 n 为任意整数 \(\geqslant 1\) ）等同，其中 \(a_{i}\) 或者是 I 的元素，或者是特殊符号 P（表示“开括号”），并且这些序列满足以下条件：
\end{enumerate}

\begin{enumerate}
\def\labelenumi{(\roman{enumi})}
\item
  序列的长度 \(n\) 等于其中出现的符号 \(\mathbf{P}\) 的个数的两倍再加 1。
\item
  对所有 \(k \leqslant n - 1\) ，符号 \(\mathbf{P}\) （出现在子序列 \((a_i)_{1 \leqslant i \leqslant k}\) 中的）的个数 \(\geqslant k / 2\) 。
\end{enumerate}

（利用《集合论》I，附录，将 M(I) 的一个元素与每个序列 \((a_i)\) 相对应；例如，序列 PPPxyzPxy 对应于字 (((xy)z)(xy))。）

¶ 13. 设 A 是一个交换环，E 是一个 A-模。A-模 \(E_{n}\) 由下式归纳地定义：

\[
\mathbf {E} _ {1} = \mathbf {E}, \quad \mathbf {E} _ {n} = \bigoplus_ {p + q = n} \left(\mathbf {E} _ {p} \otimes_ {\mathbf {A}} \mathbf {E} _ {q}\right) \quad \text { for } n \geqslant 2.
\]

我们记 \(\mathbf{ME} = \bigoplus_{n=1}^{\infty} \mathbf{E}_n\) ，并在 ME 上按下述规则定义一个 A-代数结构：对 \(x_p \in \mathbf{E}_p\) 和 \(x_q \in \mathbf{E}_q\) ， \(x_p x_q = x_p \otimes x_q \in \mathbf{E}_{p+q}\) 。

\begin{enumerate}
\def\labelenumi{(\alph{enumi})}
\item
  证明：对每个 A-代数 B 和每个 A-线性映射 \(f: \mathbf{E} \to \mathbf{B}\) ，存在唯一的代数的 A-同态 \(\mathbf{ME} \to \mathbf{B}\) 延拓 \(f\) 。ME 称为 A-模 E 的自由代数。
\item
  整数 \(a(n)\) 由下列公式对 \(n\) 归纳地定义
\end{enumerate}

\[
a (1), \quad a (n) = \sum_ {p + q = n} a (p) a (q) \quad \text { if } n \geqslant 2.
\]

证明 \(E_{n}\) 同构于 \(a(n)\) 个同构于 \(E^{\otimes n}\) 的模的直和。

\begin{enumerate}
\def\labelenumi{(\alph{enumi})}
\setcounter{enumi}{2}
\tightlist
\item
  设 \(f(\mathbf{T})\) 是形式幂级数 \(\sum_{n=1}^{\infty} a(n)\mathbf{T}^n\) 。证明
\end{enumerate}

\[
f (\mathbf {T}) = \mathbf {T} + (f (\mathbf {T})) ^ {2}.
\]

*由此推出公式

\[
f (\mathrm{T}) = \frac {1 - \sqrt {1 - 4 \mathrm{T}}}{2} \quad \text { and } \quad a (n) = 2 ^ {n - 1} \frac {1 . 3 . 5 \dots (2 n - 1)}{n !}. *
\]

\begin{enumerate}
\def\labelenumi{(\alph{enumi})}
\setcounter{enumi}{3}
\tightlist
\item
  若 I 是任意集合且 \(E = A^{(I)}\) ，证明 \(\operatorname{Lib}_{\mathbf{A}}(\mathbf{I})\) 可与 ME 等同。存在典范同态 \(\mathbf{M}(\mathbf{I}) \xrightarrow{w} \mathbf{N}^{(\mathbf{I})} \xrightarrow{l} \mathbf{N}\) ，其复合是 \(\mathbf{M}(\mathbf{I})\) 中的长度； \(\omega(m)\) 称为 \(m \in \mathbf{M}(\mathbf{I})\) 的多重次数；映射 \(\omega\) 在 \(\operatorname{Lib}_{\mathbf{A}}(\mathbf{I})\) 上定义了一个类型为 \(N^{(I)}\) 的分次。证明在此分次下，集合 \((\operatorname{Lib}_{\mathbf{A}}(\mathbf{I}))_{\alpha}\) （由多重次数为 \(\alpha \in N^{(I)}\) 的齐次元素组成）是一个自由 A-模，其秩等于
\end{enumerate}

\[
a (\alpha) = a (n) \frac {n !}{\alpha !} = 2 ^ {n - 1} \frac {1 . 3 . 5 \dots (2 n - 1)}{\alpha !}
\]

（其中 \(\alpha! = \prod_{i \in I} (\alpha(i))!\) ，而 \(n = \sum_{i \in I} \alpha(i) = l(\alpha)\) 是 \(\alpha\) 的长度。

\begin{enumerate}
\def\labelenumi{\arabic{enumi}.}
\setcounter{enumi}{13}
\tightlist
\item
  \begin{enumerate}
  \def\labelenumii{(\alph{enumii})}
  \tightlist
  \item
    设 M 为一个幺半群，其合成律记为 \(\top\) ；进一步假设 M 被一个序关系 \(x \leqslant y\) 全序化，使得关系 \(x < y\) 与 \(x' \leqslant y'\) ，或者 \(x \leqslant y\) 与 \(x' < y'\) ，蕴含 \(x \top x' < y \top y'\) 。证明：若 A 是整环，则幺半群 M 在 A 上的代数没有零因子。
  \end{enumerate}
\end{enumerate}

\begin{enumerate}
\def\labelenumi{(\alph{enumi})}
\setcounter{enumi}{1}
\item
  由此推出，如果 \(\mathbf{A}\) 是整环，则每个多项式代数 \(\mathbf{A}[(\mathbf{X}_i)_{i\in I}]\) 是整环。
\item
  若 A 是整环，且 f 和 g 是 \(\mathrm{A}[(\mathrm{X}_{i})_{i\in\mathbf{I}}]\) 中的两个多项式，使得 fg 是齐次的且 \(\neq0\) ，证明 f 和 g 都是齐次的。特别地，环 \(\mathrm{A}[(\mathrm{X}_{i})_{i\in\mathbf{I}}]\) 的可逆元素正是环 A 的可逆元素。
\end{enumerate}

\begin{enumerate}
\def\labelenumi{\arabic{enumi}.}
\setcounter{enumi}{14}
\item
  设 A 为交换环，n 为整数 \(\geqslant2\) ，使得对所有 \(a\in A\) ，方程 nx=a 有解 \(x\in A\) 。设 m 为整数 \(\geqslant1\) ，u 为 A{[}X{]} 中形如 \(X^{mn}+f(X)\) 的多项式，其中 f 的次数 \(\leqslant mn-1\) ；证明存在多项式 \(v\in A[X]\) （形如 \(X^{m}+w\) ，其中 w 的次数 \(\leqslant m-1\) ），使得 \(u-v^{n}\) 为零或次数 \(<m(n-1)\) 。
\item
  设 A 是一个交换环， \(u = \sum_{k=0}^{m} a_k X^k\) 是环 A{[}X{]} 中的一个零因子。证明：若 \(v = \sum_{k=0}^{n} b_k X^k\) 是 \(\neq 0\) 的、A{[}X{]} 中次数为 \(n\) 的元素，使得 \(uv = 0\) ，则存在多项式 \(w \neq 0\) （次数 \(\leqslant n - 1\) ），使得 \(uw = 0\) 。（化归到 \(b_0 \neq 0\) 的情形；若 \(a_k v = 0\) （对 \(0 \leqslant k \leqslant m - 1\) ），证明可取 \(w = b_0\) ；若 \(a_k v = 0\) 对于
\end{enumerate}

\[
0 \leqslant k <   p \leqslant m - 1
\]

且 \(a_{p}v \neq 0\) ，证明 \(a_{p}b_{0} = 0\) ，因此可取 \(w = \sum_{k=0}^{n-1} a_{p}b_{k+1}\mathbf{X}^{k}\) 。由此推出存在元素 \(c \neq 0\) （属于 \(\mathbf{A}\) ），使得 \(cu = 0\) 。

\begin{enumerate}
\def\labelenumi{\arabic{enumi}.}
\setcounter{enumi}{16}
\item
  设 A 是一个交换环。证明在 A 上一元多项式的代数 A{[}X{]} 中，映射 \((u,v)\mapsto u(v)\) 是一个合成律，它关于 A{[}X{]} 中的加法和乘法是结合的且左分配的。若 A 是整环，证明关系 \(u(v)=0\) 蕴含 u=0 或 v 为常数，且若 \(u\neq0\) 而 v 的次数 \textgreater0 ，则 \(u(v)\) 的次数等于 u 与 v 的次数之积。
\item
  设 K 是一个交换域，K{[}X{]} 是 K 上单不定元的多项式环。证明环 K{[}X{]} 的每个自同构 s 保持 K 不变（习题 14(c)），并在 K 上诱导出该域的一个自同构；进一步证明必然有 \(s(\mathbf{X}) = a\mathbf{X} + b\) ，其中 \(a \neq 0\) （在 K 中）， \(b \in \mathbf{K}\) （参见习题 17）。反之，证明给定一个自同构 \(\sigma\) （ \(\mathbf{K}\) 的）和两个元素 \(a, b\) （ \(\mathbf{K}\) 的，满足 \(a \neq 0\) ），就定义了且仅定义了一个自同构 \(s\) （ \(\mathbf{K}[\mathbf{X}]\) 的），使得 \(s|_{\mathbf{K}} = \sigma\) 且 \(s(\mathbf{X}) = a\mathbf{X} + b\) 。若 \(G\) 是环 \(\mathbf{K}[\mathbf{X}]\) 的所有自同构构成的群， \(N\) 是 \(G\) 的由 \(\mathbf{K}[\mathbf{X}]\) 的 K-代数结构组成的子群，证明 \(N\) 是 \(G\) 的一个正规子群，且 \(G / N\) 同构于域 \(K\) 的自同构群；群 \(N\) 同构于定义在集合 \(K^* \times K\) 上、以合成律 \((\lambda, \mu)(\lambda', \mu') = (\lambda\lambda', \lambda' \mu + \mu')\) 定义的群。
\item
  证明8个不定元中的多项式恒等式
\end{enumerate}

\[
\begin{array}{r l} & (\mathrm{X} _ {1} ^ {2} + \mathrm{X} _ {2} ^ {2} + \mathrm{X} _ {3} ^ {2} + \mathrm{X} _ {4} ^ {2}) (\mathrm{Y} _ {1} ^ {2} + \mathrm{Y} _ {2} ^ {2} + \mathrm{Y} _ {3} ^ {2} + \mathrm{Y} _ {4} ^ {2}) \\ & \qquad = (\mathrm{X} _ {1} \mathrm{Y} _ {1} - \mathrm{X} _ {2} \mathrm{Y} _ {2} - \mathrm{X} _ {3} \mathrm{Y} _ {3} - \mathrm{X} _ {4} \mathrm{Y} _ {4}) ^ {2} \\ & \qquad + (\mathrm{X} _ {1} \mathrm{Y} _ {2} + \mathrm{X} _ {2} \mathrm{Y} _ {1} + \mathrm{X} _ {3} \mathrm{Y} _ {4} - \mathrm{X} _ {4} \mathrm{Y} _ {3}) ^ {2} \\ & \qquad + (\mathrm{X} _ {1} \mathrm{Y} _ {3} + \mathrm{X} _ {3} \mathrm{Y} _ {1} + \mathrm{X} _ {4} \mathrm{Y} _ {2} - \mathrm{X} _ {2} \mathrm{Y} _ {4}) ^ {2} \\ & \qquad + (\mathrm{X} _ {1} \mathrm{Y} _ {4} + \mathrm{X} _ {4} \mathrm{Y} _ {1} + \mathrm{X} _ {2} \mathrm{Y} _ {3} - \mathrm{X} _ {3} \mathrm{Y} _ {2}) ^ {2}. \end{array}
\]

（将第 5 号的公式 (32) 应用于一个适当的四元数代数。）

\begin{enumerate}
\def\labelenumi{\arabic{enumi}.}
\setcounter{enumi}{19}
\tightlist
\item
  证明不存在关于6个不定元的多项式恒等式
\end{enumerate}

\[
(\mathrm{X} _ {1} ^ {2} + \mathrm{X} _ {2} ^ {2} + \mathrm{X} _ {3} ^ {2}) (\mathrm{Y} _ {2} ^ {2} + \mathrm{Y} _ {2} ^ {2} + \mathrm{Y} _ {3} ^ {2}) = \mathrm{P} ^ {2} + \mathrm{Q} ^ {2} + \mathrm{R} ^ {2}
\]

其中 P、Q、R 是系数在 Z 中且关于不定元 \(X_{i}\) 和 \(Y_{i}\) 的三个多项式。（注意 15 = 3·5 不能写成形式 \(p^{2} + q^{2} + r^{2}\) ，其中 p、q、r 为整数。）

\begin{enumerate}
\def\labelenumi{\arabic{enumi}.}
\setcounter{enumi}{20}
\tightlist
\item
  设 S 为具有单位元 e 的乘法幺半群，满足第 10 号中的条件 (D)，并且进一步满足以下两个条件：(1) 关系 st = e 蕴含 s = t = e；(2) 对所有 \(s \in S\) ，有限序列 \((t_{i})_{1 \leq i \leq n}\) （由不同于 e 且满足 \(t_{1}t_{2} \ldots t_{n} = s\) 的元素组成）的项数 n 小于一个仅依赖于 s 的有限数 \(v(s)\) 。
\end{enumerate}

证明：元素 \(x = \sum_{s \in S} \alpha_s s\) （属于 \(S\) 在环 \(A\) 上的全代数）可逆的必要充分条件是 \(\alpha_e\) 在 \(A\) 中可逆（化归到 \(\alpha_e = 1\) 的情形，并利用恒等式

\[
e - z ^ {n + 1} = (e - z) (e + z + \dots + z ^ {n})).
\]

\begin{enumerate}
\def\labelenumi{\arabic{enumi}.}
\setcounter{enumi}{21}
\item
  若 K 是一个交换域，证明在形式幂级数环 \(K[[X_{1}, X_{2}, \ldots, X_{p}]]\) 中只有一个极大理想，即不可逆元素的集合。另一方面，证明在多项式环 \(K[X_{1}, \ldots, X_{p}]\) 中，对 \(p \geqslant 1\) 总存在若干个不同的极大理想。
\item
  设 K 为一个交换域；证明不存在形式幂级数 \(u(\mathbf{X}, \mathbf{Y}) \in \mathbf{K}[[\mathbf{X}, \mathbf{Y}]]\) 使得，对于某个整数 \(m > 0\) ,
\end{enumerate}

\[
(\mathrm{X} + \mathrm{Y}) u (\mathrm{X}, \mathrm{Y}) = \mathrm{X} ^ {m} \mathrm{Y} ^ {m}.
\]

\begin{enumerate}
\def\labelenumi{\arabic{enumi}.}
\setcounter{enumi}{23}
\item
  设 K 是一个交换域，并设 k 是一个使得 \(k \neq 0\) （在 K 中）的整数。证明对 K{[}{[}X{]}{]} 的每个常数项等于 1 的形式幂级数 u，存在形式幂级数 \(v \in K[[X]]\) ，使得 \(v^{k} = u\) 。
\item
  设 A 为一个环（不管交换与否），并设 \(\sigma\) 为 A 的一个自同态。在加法乘积群 \(\mathbf{E} = \mathbf{A}^{\mathbf{N}}\) 上定义一个内部合成律，写作 \((\alpha_{n})(\beta_{n}) = (\gamma_{n})\) ，其中 \(\gamma_{n} = \sum_{p + q = n}\alpha_{p}\sigma^{p}(\beta_{q})\) （约定 \(\sigma^0 (\xi) = \xi\) ）。
\end{enumerate}

\begin{enumerate}
\def\labelenumi{(\alph{enumi})}
\tightlist
\item
  证明这个合成律与 E 上的加法一起，在 E 上定义了一个环结构，并以序列 \((\alpha_{n})\) （其中 \(\alpha_{0}=1\) 和 \(\alpha_{n}=0\) ，对 \(n\geqslant1\) ）为单位元 e。把每个 \(\xi\in A\) 对应到元素 \((\alpha_{n})\in E\) （满足 \(\alpha_{0}=\xi,\alpha_{n}=0\) ，对 \(n\geqslant1\) ）的映射是 A 到一个
\end{enumerate}

E 的子环，并与之等同。于是我们记 \(\sum_{n=0}^{\infty} \alpha_n X^n\) 而不是

\[
u = \sum_ {n = 0} ^ {\infty} \alpha_ {n} X ^ {n} \neq 0.
\]

\((\alpha_{n})\) ，并且 \(\mathbf{X}^{p}\beta = \sigma^{p}(\beta)\mathbf{X}^{p}\) 对所有 \(\beta \in \mathbf{A}\) 成立；若 \(u = \sum_{n=0}^{\infty} \alpha_{n}\mathbf{X}^{n} \neq 0\) ，则最小整数 \(n\) （使得 \(\alpha_{n} \neq 0\) ）称为 \(u\) 的阶，记作 \(\omega(u)\) 。

\begin{enumerate}
\def\labelenumi{(\alph{enumi})}
\setcounter{enumi}{1}
\tightlist
\item
  若 A 是一个没有零因子的环，且 \(\sigma\) 是 A 到 A 的一个子环上的同构，证明 E 是一个无零因子的环，并且
\end{enumerate}

\[
\omega (u v) = \omega (u) + \omega (v)
\]

对 \(u \neq 0\) 和 \(v \neq 0\) 。

\begin{enumerate}
\def\labelenumi{(\alph{enumi})}
\setcounter{enumi}{2}
\item
  \(u = \sum_{n=0}^{\infty} \alpha_n X^n\) 可逆的必要充分条件是 \(\alpha_0\) 在 A 中可逆。
\item
  假设 A 是一个域，且 \(\sigma\) 是 A 的一个自同构。证明 E 具有一个左分式域（I，§ 9，习题 15），并且在这个域 F 中，每个元素 \(\neq0\) 都可以唯一地写成 \(uX^{-h}\) 的形式，其中 u 是 E 中阶为 0 的元素。
\end{enumerate}

¶ *26. (a) 设 A 和 B 是 R 的两个良序子集（它们必然是可数的；参见《一般拓扑学》IV，§ 2，习题 1）。证明 A + B 是良序的，并且对所有 \(c \in \mathbf{A} + \mathbf{B}\) ，只存在有限多个有序对 \((a, b)\) （满足 \(a \in \mathbf{A}\) , \(b \in \mathbf{B}\) 和 \(a + b = c\) ）（为证明 A + B 的非空子集有最小元素，考虑它在 R 中的最大下界）。

\begin{enumerate}
\def\labelenumi{(\alph{enumi})}
\setcounter{enumi}{1}
\tightlist
\item
  设 K 为一个交换域。在向量空间 \(K^{R}\) 中，考虑由元素 \((\alpha_{x})\) 组成的向量子空间 E，其中依赖于 \((\alpha_{x})\) 的、由 \(x \in \mathbf{R}\) （满足 \(\alpha_{x} \neq 0\) ）组成的集合是良序的。对 E 中任意两个元素 \((\alpha_{x})\) , \((\beta_{x})\) ，我们记 \((\alpha_{x})(\beta_{x}) = (\gamma_{x})\) ，其中 \(\gamma_{x} = \sum_{y+z=x} \alpha_{y} \beta_{z}\) （此和式由 (a) 保证有意义）。证明这个合成律与加法一起，在 E 上定义了一个域结构；E 的元素也记作 \(\sum_{t \in \mathbb{R}} \alpha_{t} X^{t}\) ，称为系数在 K 中、具有良序实指数的形式幂级数。*
\end{enumerate}

\BourbakiExercisesSection

\begin{enumerate}
\def\labelenumi{\arabic{enumi}.}
\tightlist
\item
  设 \(\Delta\) 是一个写成加法形式的交换幺半群，其单位元记为 0，且其所有元素都是可消去的。设 E 是一个类型为 \(\Delta\) 的分次 A-代数，并假设 E 具有单位元 \(e = \sum_{\alpha \in \Delta} e_{\alpha}\) （其中 \(e_{\alpha} \in \mathbf{E}_{\alpha}\) ）。证明必然有 \(e = e_0 \in \mathbf{E}_0\) 。
\end{enumerate}
